
\Data\ID{1003019403}
\Updated{2020-09-01 11:09:31}
\Level{A}
\SubArea{Properties of functions}
\LANGen{
\Question{
Suppose each of the following tables defines function \( f \) completely. Identify which of the tables represents an odd function.}
{\(\begin{array}{|c|c|c|c|c|c|c|c|} \hline x&-5&-3& -2&0&2&3&5 \\\hline f(x) &2&-3&1&0&-1&3&-2\\ \hline\end{array}\)}{\(\begin{array}{|c|c|c|c|c|c|c|c|} \hline x& -5 & -3 & -1 & 0 & 1 & 3 & 5 \\\hline f(x) & -5 & -3 & -1 & 1 & 1 & 3 & 5 \\ \hline\end{array}\)}{\(\begin{array}{|c|c|c|c|c|c|c|c|} \hline x & -3 & -2 & -1 & 0 & 1 & 2 & 3 \\\hline f(x) & 2 & -3 &1 & 0 & 1 & -3 & 2\\ \hline\end{array}\)}{\(\begin{array}{|c|c|c|c|c|c|c|c|} \hline x & -3 & -2 & -1 & 1 & 2 & 3 & 4 \\\hline f(x) & 2 & -3 & 1 & -1 & 3 & 2 & 4\\ \hline\end{array}\)}{}{}}
\LANGpl{
\Question{
Załóżmy, że każda z poniższych tabeli określa funkcję \( f \) całkowicie. Która z przedstawionych tabeli reprezentuje funkcję nieparzystą?}
{\(\begin{array}{|c|c|c|c|c|c|c|c|} \hline x&-5&-3& -2&0&2&3&5 \\\hline f(x) &2&-3&1&0&-1&3&-2\\ \hline\end{array}\)}{\(\begin{array}{|c|c|c|c|c|c|c|c|} \hline x& -5 & -3 & -1 & 0 & 1 & 3 & 5 \\\hline f(x) & -5 & -3 & -1 & 1 & 1 & 3 & 5 \\ \hline\end{array}\)}{\(\begin{array}{|c|c|c|c|c|c|c|c|} \hline x & -3 & -2 & -1 & 0 & 1 & 2 & 3 \\\hline f(x) & 2 & -3 &1 & 0 & 1 & -3 & 2\\ \hline\end{array}\)}{\(\begin{array}{|c|c|c|c|c|c|c|c|} \hline x & -3 & -2 & -1 & 1 & 2 & 3 & 4 \\\hline f(x) & 2 & -3 & 1 & -1 & 3 & 2 & 4\\ \hline\end{array}\)}{}{}}
\LANGcs{
\Question{
Která z následujících funkcí \( f \) daných tabulkou je lichá?}
{\(\begin{array}{|c|c|c|c|c|c|c|c|} \hline x&-5&-3& -2&0&2&3&5 \\\hline f(x) &2&-3&1&0&-1&3&-2\\ \hline\end{array}\)}{\(\begin{array}{|c|c|c|c|c|c|c|c|} \hline x& -5 & -3 & -1 & 0 & 1 & 3 & 5 \\\hline f(x) & -5 & -3 & -1 & 1 & 1 & 3 & 5 \\ \hline\end{array}\)}{\(\begin{array}{|c|c|c|c|c|c|c|c|} \hline x & -3 & -2 & -1 & 0 & 1 & 2 & 3 \\\hline f(x) & 2 & -3 &1 & 0 & 1 & -3 & 2\\ \hline\end{array}\)}{\(\begin{array}{|c|c|c|c|c|c|c|c|} \hline x & -3 & -2 & -1 & 1 & 2 & 3 & 4 \\\hline f(x) & 2 & -3 & 1 & -1 & 3 & 2 & 4\\ \hline\end{array}\)}{}{}}
\LANGsk{
\Question{
Predpokladajme, že každá tabuľka určuje úplnú funkciu \( f \). Ktorá z následujúcich tabuliek určuje nepárnu funkciu?}
{\(\begin{array}{|c|c|c|c|c|c|c|c|} \hline x&-5&-3& -2&0&2&3&5 \\\hline f(x) &2&-3&1&0&-1&3&-2\\ \hline\end{array}\)}{\(\begin{array}{|c|c|c|c|c|c|c|c|} \hline x& -5 & -3 & -1 & 0 & 1 & 3 & 5 \\\hline f(x) & -5 & -3 & -1 & 1 & 1 & 3 & 5 \\ \hline\end{array}\)}{\(\begin{array}{|c|c|c|c|c|c|c|c|} \hline x & -3 & -2 & -1 & 0 & 1 & 2 & 3 \\\hline f(x) & 2 & -3 &1 & 0 & 1 & -3 & 2\\ \hline\end{array}\)}{\(\begin{array}{|c|c|c|c|c|c|c|c|} \hline x & -3 & -2 & -1 & 1 & 2 & 3 & 4 \\\hline f(x) & 2 & -3 & 1 & -1 & 3 & 2 & 4\\ \hline\end{array}\)}{}{}}
\LANGes{
\Question{
¿Cuál de la siguientes funciones \( f \) dadas por tabla es impar?}
{\(\begin{array}{|c|c|c|c|c|c|c|c|} \hline x&-5&-3& -2&0&2&3&5 \\\hline f(x) &2&-3&1&0&-1&3&-2\\ \hline\end{array}\)}{\(\begin{array}{|c|c|c|c|c|c|c|c|} \hline x& -5 & -3 & -1 & 0 & 1 & 3 & 5 \\\hline f(x) & -5 & -3 & -1 & 1 & 1 & 3 & 5 \\ \hline\end{array}\)}{\(\begin{array}{|c|c|c|c|c|c|c|c|} \hline x & -3 & -2 & -1 & 0 & 1 & 2 & 3 \\\hline f(x) & 2 & -3 &1 & 0 & 1 & -3 & 2\\ \hline\end{array}\)}{\(\begin{array}{|c|c|c|c|c|c|c|c|} \hline x & -3 & -2 & -1 & 1 & 2 & 3 & 4 \\\hline f(x) & 2 & -3 & 1 & -1 & 3 & 2 & 4\\ \hline\end{array}\)}{}{}}
\End

\Data\ID{1003019501}
\Updated{2020-12-29 05:12:30}
\Level{A}
\SubArea{Properties of functions}
\LANGen{
\Question{
Suppose function \( f \) is given completely by the next  table.
\[
\begin{array}{|c|c|c|c|c|c|c|c|} \hline x&-3&-2& -1&0&1&2&3 \\\hline f(x) &2&-3&1&0&1&-2&2\\ \hline\end{array}\]
Identify which of the following statements is true.}
{Function \( f \) has minimum at \( x= -2\) and maximum at \( x= -3\) and at \( x= 3\).}{Function \( f \) has minimum at \( x= -3\) and maximum at \( x= 2\).}{Function \( f \) has minimum at \( x= -2\) and it has no maximum.}{Function \( f \) has minimum at \( x= -3\) and maximum at \( x=3 \).}{}{}}
\LANGpl{
\Question{
Załóżmy, że funkcja  \( f \) jest określona całkowicie w tabeli:
\[
\begin{array}{|c|c|c|c|c|c|c|c|} \hline x&-3&-2& -1&0&1&2&3 \\\hline f(x) &2&-3&1&0&1&-2&2\\ \hline\end{array}\]
Które z poniższych stwierdzeń jest prawdziwe?}
{Funkcja \( f \) osiąga minimum w  \( x= -2\) i maksimum w \( x= -3\) and at \( x= 3\).}{Funkcja \( f \) osiąga minimum w \( x= -3\) i maksimum w  \( x= 2\).}{Funkcja \( f \) osiąga minimum w \( x= -2\) i nie ma maksimum.}{Funkcja \( f \)  osiąga minimum w \( x= -3\) i maksimum w \( x=3 \).}{}{}}
\LANGcs{
\Question{
Funkce \( f \) je dána tabulkou:
\[
\begin{array}{|c|c|c|c|c|c|c|c|} \hline x&-3&-2& -1&0&1&2&3 \\\hline f(x) &2&-3&1&0&1&-2&2\\ \hline\end{array}
\]
Vyberte pravdivý výrok:}
{Funkce \( f \) má v bodě \( -2\) minimum a maximum v bodech \(-3\) a \(3\).}{Funkce \( f \) má v bodě \(-3\) minimum a maximum v bodě \(2\).}{Funkce \( f \) má v bodě \(-2\) minimum a maximum nemá.}{Funkce \( f \) má v bodě \(-3\) minimum a maximum v bodě \(3\).}{}{}}
\LANGsk{
\Question{
Predpokladajme, že tabuľka určuje úplnú funkciu \( f \). 
\[
\begin{array}{|c|c|c|c|c|c|c|c|} \hline x&-3&-2& -1&0&1&2&3 \\\hline f(x) &2&-3&1&0&1&-2&2\\ \hline\end{array}
\]
 Ktoré z nasledujúcich tvrdení je pravdivé?}
{Funkcia \( f \) má v bode \( -2\) minimum a maximum v bodoch \(-3\) a \(3\).}{Funkcia \( f \) má v bode \(-3\) minimum a maximum v bode \(2\).}{Funkcia \( f \) má v bode \(-2\) minimum a maximum nemá.}{Funkcia \( f \) má v bode \(-3\) minimum a maximum v bode \(3\).}{}{}}
\LANGes{
\Question{
Supongamos que la función \( f \) viene dada por siguiente tabla:
\[
\begin{array}{|c|c|c|c|c|c|c|c|} \hline x&-3&-2& -1&0&1&2&3 \\\hline f(x) &2&-3&1&0&1&-2&2\\ \hline\end{array}\]
¿Cuál de las frases es correcta?}
{La función \( f \)tiene mínimo en \( x= -2\) y máximos en \( x= -3\) y \( x= 3\).}{La función \( f \) tiene mínimo en \( x= -3\) en máximo en \( x= 2\).}{La función \( f \) tiene mínimo en \( x= -2\) y no tiene ningún máximo.}{La función \( f \) tiene mínimo en \( x= -3\) y máximo en \( x=3 \).}{}{}}
\End

\Data\ID{1003019502}
\Updated{2020-12-29 05:12:24}
\Level{A}
\SubArea{Properties of functions}
\LANGen{
\Question{
Suppose function \( f \) is given completely by the next table. \[
\begin{array}{|c|c|c|c|c|c|c|c|} \hline x&-2&5& 9&0&-8&2&4 \\\hline f(x) &2&-3&0&-7&-1&5&4\\ \hline\end{array}\]
Identify which of the following statements is true.}
{Function \( f \) has minimum at \( x= 0 \) and maximum at \( x= 2 \).}{Function \( f \) has minimum at \( x= 0 \) and maximum at \( x= 9 \).}{Function \( f \) has minimum at \( x= -8 \) and maximum at \( x= 2 \).}{Function \( f \) has minimum at \( x= -8 \) and maximum at \( x= 9 \).}{}{}}
\LANGpl{
\Question{
Załóżmy, że funkcja  \( f \) jest całkowicie wyrażona w podanej tabeli. \[
\begin{array}{|c|c|c|c|c|c|c|c|} \hline x&-2&5& 9&0&-8&2&4 \\\hline f(x) &2&-3&0&-7&-1&5&4\\ \hline\end{array}\]
Które ze stwierdzeń jest prawdziwe?}
{Funkcja \( f \) osiąga minimum w \( x= 0 \) i maksimum w  \( x= 2 \).}{Funkcja \( f \) osiąga minimum w \( x= 0 \) i maksimum w  \( x= 9 \).}{Funkcja \( f \) osiąga minimum w \( x= -8 \) i maksimum w \( x= 2 \).}{Funkcja \( f \) osiąga minimum w \( x= -8 \) i maksimum w \( x= 9 \).}{}{}}
\LANGcs{
\Question{
Funkce \( f \) je dána tabulkou:\[
\begin{array}{|c|c|c|c|c|c|c|c|} \hline x&-2&5& 9&0&-8&2&4 \\\hline f(x) &2&-3&0&-7&-1&5&4\\ \hline\end{array}\]
Vyberte pravdivý výrok:}
{Funkce \( f \) má v bodě \(0\) minimum a maximum v bodě \(2\).}{Funkce \( f \) má v bodě \(0\) minimum a maximum v bodě \(9\).}{Funkce \( f \) má v bodě \(-8\) minimum a maximum v bodě \(2\).}{Funkce \( f \) má v bodě \(-8\) minimum a maximum v bodě \(9\).}{}{}}
\LANGsk{
\Question{
Predpokladajme, že tabuľka určuje úplnú funkciu \( f \). \[
\begin{array}{|c|c|c|c|c|c|c|c|} \hline x&-2&5& 9&0&-8&2&4 \\\hline f(x) &2&-3&0&-7&-1&5&4\\ \hline\end{array}\]
Ktoré z nasledujúcich tvrdení je pravdivé?}
{Funkcia \( f \) má v bode \(0\) minimum a maximum v bode \(2\).}{Funkcia \( f \) má v bode \(0\) minimum a maximum v bode \(9\).}{Funkcia \( f \) má v bode \(-8\) minimum a maximum v bode \(2\).}{Funkcia \( f \) má v bode \(-8\) minimum a maximum v bode \(9\).}{}{}}
\LANGes{
\Question{
Supongamos que la función \( f \) viene dada por la siguiente tabla:
 \[
\begin{array}{|c|c|c|c|c|c|c|c|} \hline x&-2&5& 9&0&-8&2&4 \\\hline f(x) &2&-3&0&-7&-1&5&4\\ \hline\end{array}\]
Identifica cuál de las declaraciones es correcta:}
{La función \( f \) tiene mínimo en \( x= 0 \) y máximo en \( x= 2 \).}{La función \( f \) tiene mínimo en \( x= 0 \) y máximo en \( x= 9 \).}{La función \( f \) tiene mínimo en \( x= -8 \) y máximo en \( x= 2 \).}{La función \( f \) tiene mínimo en \( x= -8 \) y máximo en \( x= 9 \).}{}{}}
\End

\Data\ID{1003028405}
\Updated{2020-12-30 10:12:53}
\Level{A}
\SubArea{Properties of functions}
\LANGen{
\Question{
Suppose the function \( f \) is given completely by the next  table.
\[
\begin{array}{|c|c|c|c|c|c|c|c|} \hline x&-2&-1&	0&1&2&3&4\\\hline y&0&1&0&2&3&5&4 \\\hline \end{array}
\]
Which of the following statements about the domain of the function \( f \) is true?}
{\( D(f)=\{-2; -1;0;1;2;3;4\} \)}{\( D(f)=\{0;1;2;3;4;5\} \)}{\( D(f)=\{-2;-1;0;1;2;3;4;5\} \)}{\( D(f)=[-2;4] \)}{}{}}
\LANGpl{
\Question{
Załóżmy, że funkcja  \( f \) jest podana w całości w tabeli poniżej. 
\[
\begin{array}{|c|c|c|c|c|c|c|c|} \hline x&-2&-1&	0&1&2&3&4\\\hline y&0&1&0&2&3&5&4 \\\hline \end{array}
\]
Które z poniższych stwierdzeń dotyczących dziedziny funkcji \( f \) jest prawdziwe?}
{\( D(f)=\{-2; -1;0;1;2;3;4\} \)}{\( D(f)=\{0;1;2;3;4;5\} \)}{\( D(f)=\{-2;-1;0;1;2;3;4;5\} \)}{\( D(f)=\langle-2;4\rangle \)}{}{}}
\LANGcs{
\Question{
Funkce \( f:y=f(x) \) je dána tabulkou:
\[
\begin{array}{|c|c|c|c|c|c|c|c|} \hline x&-2&-1&	0&1&2&3&4\\\hline y&0&1&0&2&3&5&4 \\\hline \end{array}
\]
Vyberte pravdivý výrok o definičním oboru funkce \( f \).}
{\( D(f)=\{-2; -1;0;1;2;3;4\} \)}{\( D(f)=\{0;1;2;3;4;5\} \)}{\( D(f)=\{-2;-1;0;1;2;3;4;5\} \)}{\( D(f)=\langle-2;4\rangle \)}{}{}}
\LANGsk{
\Question{
Predpokladajme, že tabuľka určuje úplnú funkciu \( f \). 
\[
\begin{array}{|c|c|c|c|c|c|c|c|} \hline x&-2&-1&	0&1&2&3&4\\\hline y&0&1&0&2&3&5&4 \\\hline \end{array}
\]
Ktoré tvrdenie o definičnom obore funkcie \( f \) je pravdivé?}
{\( D(f)=\{-2; -1;0;1;2;3;4\} \)}{\( D(f)=\{0;1;2;3;4;5\} \)}{\( D(f)=\{-2;-1;0;1;2;3;4;5\} \)}{\( D(f)=\langle-2;4\rangle \)}{}{}}
\LANGes{
\Question{
Supongamos que la función \( f \) viene dada por la siguiente tabla:
\[
\begin{array}{|c|c|c|c|c|c|c|c|} \hline x&-2&-1&	0&1&2&3&4\\\hline y&0&1&0&2&3&5&4 \\\hline \end{array}
\]
¿Cuál de las siguientes declaraciones sobre el dominio de la función \( f \) es correcta?}
{\( D(f)=\{-2; -1;0;1;2;3;4\} \)}{\( D(f)=\{0;1;2;3;4;5\} \)}{\( D(f)=\{-2;-1;0;1;2;3;4;5\} \)}{\( D(f)=[-2;4] \)}{}{}}
\End

\Data\ID{1003028406}
\Updated{2020-12-30 10:12:27}
\Level{A}
\SubArea{Properties of functions}
\LANGen{
\Question{
Suppose the function \( f \) is given completely by the next  table.
\[
\begin{array}{|c|c|c|c|c|c|c|c|}\hline x&-3&-2&-1&0&1&2&3 \\\hline y&-4&4&-4&4&-4&4&-4 \\\hline \end{array}
\]
Which of the following statements about the range of the function \( f \) is true?}
{\( H(f)=\{-4; 4\} \)}{\( H(f)=\{-3;-2;-1;0;1;2;3\} \)}{\( H(f)=[-4;4] \)}{\( H(f)=(-4;4) \)}{}{}}
\LANGpl{
\Question{
Załóżmy, ze funkcja  \( f \) jest całkowicie wyrażona w tabeli poniżej.
\[
\begin{array}{|c|c|c|c|c|c|c|c|}\hline x&-3&-2&-1&0&1&2&3 \\\hline y&-4&4&-4&4&-4&4&-4 \\\hline \end{array}
\]
Które ze stwierdzeń dotyczących zakresu funkcji \( f \)  jest prawdziwe?}
{\( H(f)=\{-4; 4\} \)}{\( H(f)=\{-3;-2;-1;0;1;2;3\} \)}{\( H(f)=\langle-4;4\rangle \)}{\( H(f)=(-4;4) \)}{}{}}
\LANGcs{
\Question{
Funkce  \( f:y=f(x) \) je dána tabulkou:
\[
\begin{array}{|c|c|c|c|c|c|c|c|}\hline x&-3&-2&-1&0&1&2&3 \\\hline y&-4&4&-4&4&-4&4&-4 \\\hline \end{array}
\]
Vyberte pravdivý výrok o oboru hodnot funkce \( f \).}
{\( H(f)=\{-4; 4\} \)}{\( H(f)=\{-3;-2;-1;0;1;2;3\} \)}{\( H(f)=\langle-4;4\rangle \)}{\( H(f)=(-4;4) \)}{}{}}
\LANGsk{
\Question{
Predpokladajme, že tabuľka určuje úplnú funkciu \( f \). 
\[
\begin{array}{|c|c|c|c|c|c|c|c|}\hline x&-3&-2&-1&0&1&2&3 \\\hline y&-4&4&-4&4&-4&4&-4 \\\hline \end{array}
\]
Ktoré tvrdenie o obore hodnôt funkcie \( f \) je pravdivé?}
{\( H(f)=\{-4; 4\} \)}{\( H(f)=\{-3;-2;-1;0;1;2;3\} \)}{\( H(f)=\langle-4;4\rangle \)}{\( H(f)=(-4;4) \)}{}{}}
\LANGes{
\Question{
Supongamos que la función \( f \) viene dada por la siguiente tabla:
\[
\begin{array}{|c|c|c|c|c|c|c|c|}\hline x&-3&-2&-1&0&1&2&3 \\\hline y&-4&4&-4&4&-4&4&-4 \\\hline \end{array}
\]
¿Cuál de las siguientes declaraciones sobre el rango de la función \( f \) es correcta?}
{\( H(f)=\{-4; 4\} \)}{\( H(f)=\{-3;-2;-1;0;1;2;3\} \)}{\( H(f)=[-4;4] \)}{\( H(f)=(-4;4) \)}{}{}}
\End

\Data\ID{1003030801}
\Updated{2020-09-01 02:09:14}
\Level{A}
\SubArea{Properties of functions}
\LANGen{
\Question{
Suppose each of the following tables defines functions completely. Identify which of the tables represents the decreasing function.}
{\(\begin{array}{|c|c|c|c|c|c|c|c|}
\hline x &-1 & -2 & 0 & -3 & 3 & 2 & 1 \\\hline f(x) & 3&4&-1&5&-5&-4&-3 \\\hline
\end{array}\)}{\(\begin{array}{|c|c|c|c|c|c|c|c|}
\hline x &3 & 2 & 1 & 0 & -1 & -2 & -3 \\\hline h(x) & 5&4&3&2&0&-1&-2 \\\hline
\end{array}
\)}{\(\begin{array}{|c|c|c|c|c|c|c|c|}
\hline x &-3 & -2 & -1 & 0 & 1 & 2 & 3 \\\hline g(x) & 3&2&1&0&3&2&1 \\\hline
\end{array}
\)}{\(\begin{array}{|c|c|c|c|c|c|c|c|}
\hline x &-1 & -2 & 0 & -3 & 3 & 2 & 1 \\\hline m(x) & 3&4&-3&5&-5&-4&-3 \\\hline
\end{array}\)}{}{}}
\LANGpl{
\Question{
Załóżmy, że każda z poniższych tabeli określa funkcję całkowicie. Która z tabeli przedstawia funkcję malejącą?}
{\(\begin{array}{|c|c|c|c|c|c|c|c|}
\hline x &-1 & -2 & 0 & -3 & 3 & 2 & 1 \\\hline f(x) & 3&4&-1&5&-5&-4&-3 \\\hline
\end{array}\)}{\(\begin{array}{|c|c|c|c|c|c|c|c|}
\hline x &3 & 2 & 1 & 0 & -1 & -2 & -3 \\\hline h(x) & 5&4&3&2&0&-1&-2 \\\hline
\end{array}
\)}{\(\begin{array}{|c|c|c|c|c|c|c|c|}
\hline x &-3 & -2 & -1 & 0 & 1 & 2 & 3 \\\hline g(x) & 3&2&1&0&3&2&1 \\\hline
\end{array}
\)}{\(\begin{array}{|c|c|c|c|c|c|c|c|}
\hline x &-1 & -2 & 0 & -3 & 3 & 2 & 1 \\\hline m(x) & 3&4&-3&5&-5&-4&-3 \\\hline
\end{array}\)}{}{}}
\LANGcs{
\Question{
Která z následujících funkcí zadaných tabulkou je klesající?}
{\(\begin{array}{|c|c|c|c|c|c|c|c|}
\hline x &-1 & -2 & 0 & -3 & 3 & 2 & 1 \\\hline f(x) & 3&4&-1&5&-5&-4&-3 \\\hline
\end{array}
\)}{\(\begin{array}{|c|c|c|c|c|c|c|c|}
\hline x &3 & 2 & 1 & 0 & -1 & -2 & -3 \\\hline h(x) & 5&4&3&2&0&-1&-2 \\\hline
\end{array}
\)}{\(\begin{array}{|c|c|c|c|c|c|c|c|}
\hline x &-3 & -2 & -1 & 0 & 1 & 2 & 3 \\\hline g(x) & 3&2&1&0&3&2&1 \\\hline
\end{array}\)}{\(\begin{array}{|c|c|c|c|c|c|c|c|}
\hline x &-1 & -2 & 0 & -3 & 3 & 2 & 1 \\\hline m(x) & 3&4&-3&5&-5&-4&-3 \\\hline
\end{array}\)}{}{}}
\LANGsk{
\Question{
Predpokladajme, že tabuľka určuje úplnú funkciu \( f \). Ktorá z následujúcich tabuliek určuje klesajúcu funkciu?}
{\(\begin{array}{|c|c|c|c|c|c|c|c|}
\hline x &-1 & -2 & 0 & -3 & 3 & 2 & 1 \\\hline f(x) & 3&4&-1&5&-5&-4&-3 \\\hline
\end{array}
\)}{\(\begin{array}{|c|c|c|c|c|c|c|c|}
\hline x &3 & 2 & 1 & 0 & -1 & -2 & -3 \\\hline h(x) & 5&4&3&2&0&-1&-2 \\\hline
\end{array}
\)}{\(\begin{array}{|c|c|c|c|c|c|c|c|}
\hline x &-3 & -2 & -1 & 0 & 1 & 2 & 3 \\\hline g(x) & 3&2&1&0&3&2&1 \\\hline
\end{array}\)}{\(\begin{array}{|c|c|c|c|c|c|c|c|}
\hline x &-1 & -2 & 0 & -3 & 3 & 2 & 1 \\\hline m(x) & 3&4&-3&5&-5&-4&-3 \\\hline
\end{array}\)}{}{}}
\LANGes{
\Question{
¿Cuál de las tablas siguientes representa una función decreciente?}
{\(\begin{array}{|c|c|c|c|c|c|c|c|}
\hline x &-1 & -2 & 0 & -3 & 3 & 2 & 1 \\\hline f(x) & 3&4&-1&5&-5&-4&-3 \\\hline
\end{array}\)}{\(\begin{array}{|c|c|c|c|c|c|c|c|}
\hline x &3 & 2 & 1 & 0 & -1 & -2 & -3 \\\hline h(x) & 5&4&3&2&0&-1&-2 \\\hline
\end{array}
\)}{\(\begin{array}{|c|c|c|c|c|c|c|c|}
\hline x &-3 & -2 & -1 & 0 & 1 & 2 & 3 \\\hline g(x) & 3&2&1&0&3&2&1 \\\hline
\end{array}
\)}{\(\begin{array}{|c|c|c|c|c|c|c|c|}
\hline x &-1 & -2 & 0 & -3 & 3 & 2 & 1 \\\hline m(x) & 3&4&-3&5&-5&-4&-3 \\\hline
\end{array}\)}{}{}}
\End

\Data\ID{1003030906}
\Updated{2020-07-15 03:07:12}
\Level{A}
\SubArea{Properties of functions}
\LANGen{
\Question{
The function \( f \) has a maximum and has not a minimum. Identify which of the following statements is false, i.e. find the statement that is false for at least one of the functions that meet the given conditions.}
{The function \( f \) is not bounded above.}{The function \( f \) is bounded below.}{The function \( f \) is not bounded below.}{The function \( f \) is not bounded.}{}{}}
\LANGpl{
\Question{
Funkcja \( f \) posiada maksimum, ale nie ma minimum. Określ, które z poniższych stwierdzeń jest fałszywe, tzn. wybierz stwierdzenie, które jest fałszywe przynajmniej dla jednej funkcji, która spełnia poniższe warunki.}
{Funkcja  \( f \) nie jest ograniczona z góry.}{Funkcja  \( f \) jest ograniczona z dołu.}{Funkcja \( f \) nie jest ograniczona z dołu.}{Funkcja \( f \) nie jest ograniczona.}{}{}}
\LANGcs{
\Question{
Funkce \( f \) má maximum a nemá minimum. Vyberte nepravdivý výrok, tj. výrok, který neplatí alespoň pro jednu z funkcí splňujících dané podmínky.}
{Funkce  \( f \) není shora omezená.}{Funkce  \( f \) je zdola omezená.}{Funkce  \( f \) není zdola omezená.}{Funkce  \( f \) není omezená.}{}{}}
\LANGsk{
\Question{
Funkcia \( f \) má maximum a nemá minimum. Vyberte nepravdivý výrok, t.j. výrok, ktorý neplatí aspoň pre jednu z funkcií spĺňajúcich dané podmienky.}
{Funkcia \( f \) nie je zhora ohraničená.}{Funkcia \( f \) je zdola ohraničená.}{Funkcia \( f \) nie je zdola ohraničená.}{Funkcia \( f \) nie je ohraničená.}{}{}}
\LANGes{
\Question{
The function \( f \) has a maximum and has not a minimum. Identify which of the following statements is false, i.e. find the statement that is false for at least one of the functions that meet the given conditions.}
{The function \( f \) is not bounded above.}{The function \( f \) is bounded below.}{The function \( f \) is not bounded below.}{The function \( f \) is not bounded.}{}{}}
\End

\Data\ID{1003048501}
\Updated{2020-12-30 10:12:15}
\Level{A}
\SubArea{Properties of functions}
\LANGen{
\Question{
Let \( f \) be a periodic function with period \( 5 \). The table below shows some input and corresponding output values of \( f \).
\[\begin{array}{|c|c|c|c|c|c|c|c|} \hline x & -1.5 & -1 & 0 & 1 & 2 & 3 & 4 \\\hline f(x) & 0 & 4 & 1 & -1 & 3 & 2 & 4 \\\hline \end{array}\]
Identify which of the following statements is false.}
{\( f(-12)=3 \)}{\( f(5)=1 \)}{\( f(12)=3 \)}{\( f(3.5)=0 \)}{}{}}
\LANGpl{
\Question{
Niech  \( f \) będzie funkcją okresową z okresem  \( 5 \).  Poniższa tabela pokazuje niektóre dane wejściowe i odpowiadające im wartości wyjściowe \( f \).
\[\begin{array}{|c|c|c|c|c|c|c|c|} \hline x & -1{,}5 & -1 & 0 & 1 & 2 & 3 & 4 \\\hline f(x) & 0 & 4 & 1 & -1 & 3 & 2 & 4 \\\hline \end{array}\]
Określ, które z poniższych stwierdzeń jest  fałszywe.}
{\( f(-12)=3 \)}{\( f(5)=1 \)}{\( f(12)=3 \)}{\( f(3{,}5)=0 \)}{}{}}
\LANGcs{
\Question{
Funkce \( f \) je periodická s periodou  \( 5 \). Některé její funkční hodnoty jsou uvedeny v následující tabulce.
\[\begin{array}{|c|c|c|c|c|c|c|c|} \hline x & -1{,}5 & -1 & 0 & 1 & 2 & 3 & 4 \\\hline f(x) & 0 & 4 & 1 & -1 & 3 & 2 & 4 \\\hline \end{array}\]
Vyberte nepravdivý výrok.}
{\( f(-12)=3 \)}{\( f(5)=1 \)}{\( f(12)=3 \)}{\( f(3{,}5)=0 \)}{}{}}
\LANGsk{
\Question{
Daná je periodická funkcia \( f \) s periódou \( 5 \). Tabuľka ukazuje niektoré jej funkčné hodnoty.
\[\begin{array}{|c|c|c|c|c|c|c|c|} \hline x & -1{,}5 & -1 & 0 & 1 & 2 & 3 & 4 \\\hline f(x) & 0 & 4 & 1 & -1 & 3 & 2 & 4 \\\hline \end{array}\]
Vyberte tvrdenie, ktoré je nepravdivé.}
{\( f(-12)=3 \)}{\( f(5)=1 \)}{\( f(12)=3 \)}{\( f(3{,}5)=0 \)}{}{}}
\LANGes{
\Question{
Sea \( f \) una función periódica de periodo \( 5 \). En la tabla hay algunas entradas y los valores de \( f \) correspondientes.
\[\begin{array}{|c|c|c|c|c|c|c|c|} \hline x & -1.5 & -1 & 0 & 1 & 2 & 3 & 4 \\\hline f(x) & 0 & 4 & 1 & -1 & 3 & 2 & 4 \\\hline \end{array}\]
Identifica cuál de las afirmaciones es falsa.}
{\( f(-12)=3 \)}{\( f(5)=1 \)}{\( f(12)=3 \)}{\( f(3.5)=0 \)}{}{}}
\End

\Data\ID{1103019404}
\Updated{2020-09-08 10:09:00}
\Level{A}
\SubArea{Properties of functions}
\LANGen{
\Question{
Identify which of the given graphs is an even function graph.}
{}{}{}{}{}{}}
\LANGpl{
\Question{
Który z poniższych wykresów przedstawia funkcję parzystą?}
{}{}{}{}{}{}}
\LANGcs{
\Question{
Na kterém obrázku je graf sudé funkce?}
{}{}{}{}{}{}}
\LANGsk{
\Question{
Na ktorom obrázku je graf párnej funkcie?}
{}{}{}{}{}{}}
\LANGes{
\Question{
NA}
{}{}{}{}{}{}}

\IMGanswer{1003019404a_0.pdf}
\IMGanswer{1003019404b_0.pdf}
\IMGanswer{1003019404c_0.pdf}
\IMGanswer{1003019404d_0.pdf}\End

\Data\ID{1103019503}
\Updated{2020-12-30 10:12:13}
\Level{A}
\SubArea{Properties of functions}
\IMG{1003019503_0.pdf}
\LANGen{
\Question{
Function \( f \) is given by the graph. Identify which of the following statements is true.}
{Function \( f \) has minimum at \( x=0 \) and maximum at \( x=5 \).}{Function \( f \) has minimum at \( x=-5 \) and maximum at \( x=5 \).}{Function \( f \) has minimum at \( x=-1 \) and maximum at \( x=4 \).}{Function \( f \) has neither minimum nor maximum.}{}{}}
\LANGpl{
\Question{
Wykres funkcji \( f \) przedstawiono poniżej. Które z podanych stwierdzeń jest prawdziwe?}
{Funkcja \( f \) osiąga minimum w \( x=0 \) i maksimum w  \( x=5 \).}{Funkcja \( f \) osiąga minimum w \( x=-5 \) i maksimum w  \( x=5 \).}{Funkcja  \( f \) osiąga minimum w  \( x=-1 \) i maksimum w \( x=4 \).}{Funkcja  \( f \) nie ma ani maksimum ani minimum.}{}{}}
\LANGcs{
\Question{
Funkce \( f \) je dána grafem. Vyberte pravdivý výrok:}
{Funkce \( f \) má v bodě  \(0\) minimum a maximum v bodě \(5\).}{Funkce \( f \) má v bodě  \(-5\) minimum a maximum v bodě \(5\).}{Funkce \( f \) má v bodě  \(-1\) minimum a maximum v bodě \(4\).}{Funkce \( f \) nemá minimum ani maximum.}{}{}}
\LANGsk{
\Question{
Na obrázku je graf funkcie \( f \). Ktoré z nasledujúcich tvrdení o funkcii \( f \) je pravdivé?}
{Funkcia \( f \) má v bode  \(0\) minimum a maximum v bode \(5\).}{Funkcia \( f \) má v bode  \(-5\) minimum a maximum v bode \(5\).}{Funkcia \( f \) má v bode  \(-1\) minimum a maximum v bode \(4\).}{Funkcia \( f \) nemá minimum ani maximum.}{}{}}
\LANGes{
\Question{
La función \( f \) viene dada por la gráfica. ¿Cuál de las declaraciones es verdadera?}
{La función \( f \) tiene mínimo en \( x=0 \) y máximo en \( x=5 \).}{La función \( f \) tiene mínimo en \( x=-5 \) y máximo en \( x=5 \).}{La función \( f \) tiene mínimo en \( x=-1 \) y máximo en \( x=4 \).}{La función \( f \) no tiene mínimo ni máximo.}{}{}}
\End

\Data\ID{1103019504}
\Updated{2020-09-08 11:09:36}
\Level{A}
\SubArea{Properties of functions}
\LANGen{
\Question{
Identify which of the graphs represents a function without a maximum.}
{}{}{}{}{}{}}
\LANGpl{
\Question{
Który z poniższych wykresów przedstawia funkcję nie posiadającą maksimum?}
{}{}{}{}{}{}}
\LANGcs{
\Question{
Na kterém obrázku je graf funkce, která nemá maximum?}
{}{}{}{}{}{}}
\LANGsk{
\Question{
Na ktorom obrázku je graf funkcie, ktorá nemá maximum?}
{}{}{}{}{}{}}
\LANGes{
\Question{
NA}
{}{}{}{}{}{}}

\IMGanswer{1003019504a.pdf}
\IMGanswer{1003019504b.pdf}
\IMGanswer{1003019504c.pdf}
\IMGanswer{1003019504d.pdf}\End

\Data\ID{1103025601}
\Updated{2020-12-30 10:12:54}
\Level{A}
\SubArea{Properties of functions}
\IMG{1103025601.pdf}
\LANGen{
\Question{
The function \( f \) is given by the graph. Which of the following statements is true?}
{The function \( f \) has a  minimum and a  maximum at every \( x \) of its domain.}{The function \( f \) has the minimum at \( x=-6 \)  and the maximum at \( x=3 \).}{The function \( f \) has  the maximum at  \( x=3 \)  and it has no minimum.}{The function \( f \) has neither  a minimum nor  a maximum.}{}{}}
\LANGpl{
\Question{
Funkcja  \( f \) przedstawiona jest za pomocą następującego wykresu. Które z poniższych zdań jest zgodne z prawdą?}
{Funkcja \( f \) ma minimum i maksimum w każdym \( x \) swojej dziedziny.}{Funkcja \( f \) ma minimum w \( x=-6 \) i maksimum w \( x=3 \).}{Funkcja  \( f \) osiąga maksimum w   \( x=3 \)  i nie posiada minimum.}{Funkcja \( f \) nie posiada minimum ani maksimum.}{}{}}
\LANGcs{
\Question{
Funkce \( f:y=f(x) \) je dána grafem. Vyberte pravdivý výrok.}
{Funkce \( f \) má v každém bodě svého definičního oboru minimum i maximum.}{Funkce \( f \) má v bodě \( x=-6 \) minimum a v bodě \( x=3 \) maximum.}{Funkce \( f \) má maximum v bodě \( x=3 \) a minimum nemá.}{Funkce \( f \) nemá ani minimum ani maximum.}{}{}}
\LANGsk{
\Question{
Na obrázku je graf funkcie \( f \). Ktoré z nasledujúcich tvrdení o funkcii \( f \) je pravdivé?}
{Funkcia \( f \) má v každom bode svojho definičného oboru minimum aj maximum.}{Funkcia \( f \) má v bode \( x=-6 \) minimum a v bode \( x=3 \) maximum.}{Funkcia \( f \) má maximum v bode \( x=3 \) a minimum nemá.}{Funkcia \( f \) nemá ani minimum ani maximum.}{}{}}
\LANGes{
\Question{
La función \( f \) viene dada por la gráfica. ¿Cuál declaración es correcta?}
{La función \( f \) tiene mínimo y máximo en cada \( x \) de su dominio.}{La función \( f \) tiene mínimo en \( x=-6 \)  y máximo en \( x=3 \).}{La función \( f \) tiene máximo en \( x=3 \)  y no tiene mínimo.}{La función \( f \) no tiene mínimo ni máximo.}{}{}}
\End

\Data\ID{1103025602}
\Updated{2020-09-08 11:09:41}
\Level{A}
\SubArea{Properties of functions}
\LANGen{
\Question{
Identify which of the graphs represents the function with the minimum at \( x=3 \).}
{}{}{}{}{}{}}
\LANGpl{
\Question{
Który z poniższych wykresów przedstawia funkcję o minimum w punkcie \( x=3 \)?}
{}{}{}{}{}{}}
\LANGcs{
\Question{
Najděte graf funkce, která má minimum v bodě  \( x = 3 \).}
{}{}{}{}{}{}}
\LANGsk{
\Question{
Vyberte graf funkcie, ktorá má minimum v \( x=3 \).}
{}{}{}{}{}{}}
\LANGes{
\Question{
NA}
{}{}{}{}{}{}}

\IMGanswer{1103025602a_0.pdf}
\IMGanswer{1103025602b_0.pdf}
\IMGanswer{1103025602c_0.pdf}
\IMGanswer{1103025602d_0.pdf}\End

\Data\ID{1103028408}
\Updated{2020-12-30 10:12:40}
\Level{A}
\SubArea{Properties of functions}
\IMG{1103028408.pdf}
\LANGen{
\Question{
The function \( f \) is given by the graph.
Which of the statements about the domain and the range of the function \( f \) is true?}
{\( D(f) =(-2;3]; H(f)= (-1;3] \)}{\( D(f) =(-1;3] ; H(f)=(-2;3] \)}{\( D(f) =(-2;3] ; H(f)=(-1;1.5] \)}{\( D(f) =[-2;3] ; H(f)=[-1;3] \)}{}{}}
\LANGpl{
\Question{
Wykres funkcji  \( f \) przedstawiono poniżej. 
Które ze stwierdzeń dotyczących dziedziny i zakresu funkcji  \( f \) jest prawdziwe?}
{\( D(f) =(-2;3\rangle; H(f)= (-1;3\rangle \)}{\( D(f) =(-1;3\rangle; H(f)=(-2;3\rangle \)}{\( D(f) =(-2;3\rangle; H(f)=(-1;1{,}5\rangle \)}{\( D(f) =\langle-2;3\rangle; H(f)=\langle-1;3\rangle \)}{}{}}
\LANGcs{
\Question{
Funkce \( f:y=f(x) \) je dána grafem.
Vyberte pravdivý výrok o definičním oboru \(D(f)\) a oboru hodnot \(H(f)\) funkce f.}
{\( D(f) =(-2;3\rangle; H(f)= (-1;3\rangle \)}{\( D(f) =(-1;3\rangle; H(f)=(-2;3\rangle \)}{\( D(f) =(-2;3\rangle; H(f)=(-1;1{,}5\rangle \)}{\( D(f) =\langle-2;3\rangle; H(f)=\langle-1;3\rangle \)}{}{}}
\LANGsk{
\Question{
Na obrázku je graf funkcie \( f \).
Ktoré tvrdenie o definičnom obore a obore hodnôt funkcie \( f \) je pravdivé?}
{\( D(f) =(-2;3\rangle; H(f)= (-1;3\rangle \)}{\( D(f) =(-1;3\rangle; H(f)=(-2;3\rangle \)}{\( D(f) =(-2;3\rangle; H(f)=(-1;1{,}5\rangle \)}{\( D(f) =\langle-2;3\rangle; H(f)=\langle-1;3\rangle \)}{}{}}
\LANGes{
\Question{
La función \( f \) viene dada por la gráfica.
¿Cuál declaración sobre el dominio y el rango de la función \( f \) es verdadera?}
{\( D(f) =(-2;3]; H(f)= (-1;3] \)}{\( D(f) =(-1;3] ; H(f)=(-2;3] \)}{\( D(f) =(-2;3] ; H(f)=(-1;1.5] \)}{\( D(f) =[-2;3] ; H(f)=[-1;3] \)}{}{}}
\End

\Data\ID{1103028409}
\Updated{2020-12-30 10:12:19}
\Level{A}
\SubArea{Properties of functions}
\IMG{1103028409.pdf}
\LANGen{
\Question{
The function \( f \) is given by the graph. Which of the statements about the domain and the range of the function \( f \) is true?}
{\( D(f) =[-3;4]; H(f)=[-2;2)\cup(2; 3]\cup\{5\} \)}{\( D(f) =[-3;1)\cup(1; 4]; H(f)=[-2; 2)\cup(2; 3] \)}{\( D(f)=[-3;4]; H(f)=[-2;5] \)}{\( D(f) =[-3;4]; H(f)=[-2;3]\cup\{5\} \)}{}{}}
\LANGpl{
\Question{
Wykres funkcji \( f \) przedstawiono poniżej.  Które ze stwierdzeń dotyczących dziedziny i zakresu funkcji  \( f \) jest prawdziwe?}
{\( D(f) =\langle-3;4\rangle; H(f)=\langle-2;2)\cup(2; 3\rangle\cup\{5\} \)}{\( D(f) =\langle-3;1)\cup(1; 4\rangle; H(f)=\langle-2; 2)\cup(2; 3\rangle \)}{\( D(f)=\langle-3;4\rangle; H(f)=\langle-2;5\rangle \)}{\( D(f) =\langle-3;4\rangle; H(f)=\langle-2;3\rangle\cup\{5\} \)}{}{}}
\LANGcs{
\Question{
Funkce \( f:y=f(x) \) je dána grafem. Vyberte pravdivý výrok o definičním oboru  \(D(f)\)  a oboru hodnot \(H(f)\) funkce f.}
{\( D(f) =\langle-3;4\rangle ;H(f)=\langle-2;2)\cup(2; 3\rangle\cup\{5\} \)}{\( D(f) =\langle-3;1)\cup(1; 4\rangle ; H(f)=\langle-2; 2)\cup(2; 3\rangle \)}{\( D(f)=\langle-3;4\rangle ;H(f)=\langle-2;5\rangle \)}{\( D(f) =\langle-3;4\rangle;H(f)=\langle-2;3\rangle\cup\{5\} \)}{}{}}
\LANGsk{
\Question{
Na obrázku je graf funkcie \( f \).
Ktoré tvrdenie o definičnom obore a obore hodnôt funkcie \( f \) je pravdivé?}
{\( D(f) =\langle-3;4\rangle; H(f)=\langle-2;2)\cup(2; 3\rangle\cup\{5\} \)}{\( D(f) =\langle-3;1)\cup(1; 4\rangle; H(f)=\langle-2; 2)\cup(2; 3\rangle \)}{\( D(f)=\langle-3;4\rangle; H(f)=\langle-2;5\rangle \)}{\( D(f) =\langle-3;4\rangle; H(f)=\langle-2;3\rangle\cup\{5\} \)}{}{}}
\LANGes{
\Question{
La función \( f \) viene dada por la gráfica. ¿Cuál declaración sobre el Dominio y el Rango de la función \( f \) es verdadera?}
{\( D(f) =[-3;4]; R(f)=[-2;2)\cup(2; 3]\cup\{5\} \)}{\( D(f) =[-3;1)\cup(1; 4]; R(f)=[-2; 2)\cup(2; 3] \)}{\( D(f)=[-3;4]; R(f)=[-2;5] \)}{\( D(f) =[-3;4]; R(f)=[-2;3]\cup\{5\} \)}{}{}}
\End

\Data\ID{1103028410}
\Updated{2020-12-30 10:12:45}
\Level{A}
\SubArea{Properties of functions}
\LANGen{
\Question{
Identify which of the graphs represents the function \( f \) with  the domain \( [-4; 2) \) and the range  \( (-5; 3] \).}
{}{}{}{}{}{}}
\LANGpl{
\Question{
Który z poniższych wykresów przedstawia funkcję \( f \) z dziedziną  \( \langle-4; 2) \) i zakresem \( (-5; 3\rangle \)?}
{}{}{}{}{}{}}
\LANGcs{
\Question{
Určete graf funkce, jejímž definičním oborem je interval  \( \langle-4; 2) \) a oborem hodnot interval \( (-5; 3\rangle \).}
{}{}{}{}{}{}}
\LANGsk{
\Question{
Určte graf funkcie, ktorej definičný obor je interval \( \langle-4; 2) \) a obor hodnôt je interval \( (-5; 3\rangle \).}
{}{}{}{}{}{}}
\LANGes{
\Question{
NA}
{}{}{}{}{}{}}

\IMGanswer{1103028410a_0.pdf}
\IMGanswer{1103028410b_0.pdf}
\IMGanswer{1103028410c_0.pdf}
\IMGanswer{1103028410d_0.pdf}\End

\Data\ID{1103030802}
\Updated{2020-12-30 10:12:35}
\Level{A}
\SubArea{Properties of functions}
\IMG{1103030802.pdf}
\LANGen{
\Question{
The function \( f \) is given by the graph. Identify which of the following statements is true.}
{The function \( f \) is neither increasing nor decreasing.}{The function \( f \) is increasing.}{The function \( f \) is non-decreasing.}{The function \( f \) is increasing in the interval \( [ -4;1] \).}{}{}}
\LANGpl{
\Question{
Funkcję  \( f \) przedstawiono za pomocą poniższego wykresu. Które z podanych stwierdzeń jest prawdziwe?}
{Funkcja  \( f \) nie jest ani rosnąca ani malejąca.}{Funkcja  \( f \) jest rosnąca.}{Funkcja  \( f \) jest niemalejąca.}{Funkcja  \( f \) jest rosnąca w przedziale \( \langle -4;1\rangle \).}{}{}}
\LANGcs{
\Question{
Funkce \( f \) je dána grafem. Vyberte pravdivý výrok.}
{Funkce \( f \) není ani rostoucí ani klesající.}{Funkce \( f \) je rostoucí.}{Funkce \( f \) je neklesající.}{Funkce \( f \) je rostoucí v intervalu \( \langle -4;1\rangle \).}{}{}}
\LANGsk{
\Question{
Na obrázku je graf funkcie \( f \). Ktoré z nasledujúcich tvrdení o funkcii \( f \) je pravdivé?}
{Funkcia \( f \) nie je rastúca ani klesajúca.}{Funkcia \( f \) je rastúca.}{Funkcia \( f \) je neklesajúca.}{Funkcia \( f \) je rastúca na intervalu \( \langle -4;1\rangle \).}{}{}}
\LANGes{
\Question{
La función \( f \) viene dada por la gráfica. ¿Cuál de las declaraciones siguientes es verdadera?}
{La función \( f \) no es decreciente ni creciente.}{La función \( f \) es creciente.}{La función \( f \) es no-decreciente.}{La función \( f \) es creciente en el intervalo \( [ -4;1] \).}{}{}}
\End

\Data\ID{1103030803}
\Updated{2020-12-30 10:12:10}
\Level{A}
\SubArea{Properties of functions}
\IMG{1103030803.pdf}
\LANGen{
\Question{
There is a part of the graph of the function \( f(x)=x^3 \) in the picture. Identify which of the following statements is true.}
{The function \( f \) is increasing in the interval \( [ -1;1 ] \).}{The function \( f \) is decreasing in the interval  \( [ -1;1 ] \).}{The function \( f \) is non-decreasing and it is not increasing in the interval \( [ -1;1 ] \).}{The function \( f \) is non-increasing in the interval \( [ -1;1 ] \).}{}{}}
\LANGpl{
\Question{
Na rysunku przedstawiono część wykresu funkcji \( f(x)=x^3 \). Które z podanych zdań jest prawdziwe?}
{Funkcja  \( f \) jest rosnąca w przedziale \( \langle -1;1 \rangle \).}{Funkcja  \( f \) jest malejąca w przedziale \( \langle -1;1 \rangle \).}{Funkcja  \( f \) funkcja jest niemalejąca i nie jest rosnąca w przedziale \( \langle -1;1 \rangle \).}{Funkcja  \( f \) jest nierosnąca w przedziale  \( \langle -1;1 \rangle \).}{}{}}
\LANGcs{
\Question{
Na obrázku je část grafu funkce \( f(x)=x^3 \). Vyberte pravdivý výrok.}
{Funkce \( f \) je v intervalu \( \langle -1;1 \rangle \) rostoucí.}{Funkce \( f \) je v intervalu \( \langle -1;1 \rangle \) klesající.}{Funkce \( f \) je v intervalu \( \langle -1;1 \rangle \) neklesající a zároveň není v tomto intervalu rostoucí.}{Funkce \( f \) je v intervalu \( \langle -1;1 \rangle \) nerostoucí.}{}{}}
\LANGsk{
\Question{
Na obrázku je časť grafu funkcie \( f(x)=x^3 \). Ktoré z nasledujúcich tvrdení o funkcii \( f \)  je pravdivé?}
{Funkcia \( f \) je na intervale \( \langle -1;1 \rangle \) rastúca.}{Funkcia \( f \) je na intervale \( \langle -1;1 \rangle \) klesajúca.}{Funkcia \( f \) je na intervale \( \langle -1;1 \rangle \) neklesajúca a zároveň nie je v tomto intervale rastúca.}{Funkcia \( f \) je na intervale \( \langle -1;1 \rangle \) nerastúca.}{}{}}
\LANGes{
\Question{
En el dibujo hay una parte de la gráfica de la función \( f(x)=x^3 \). ¿Cuál de las siguientes declaraciones es verdadera?}
{La función \( f \) es creciente en el intervalo \( [ -1;1 ] \).}{La función \( f \) es decreciente en el intervalo  \( [ -1;1 ] \).}{La función \( f \) es no-decreciente y no es creciente en el intervalo \( [ -1;1 ] \).}{La función \( f \) es no-creciente en el intervalo \( [ -1;1 ] \).}{}{}}
\End

\Data\ID{1103030806}
\Updated{2020-12-30 10:12:31}
\Level{A}
\SubArea{Properties of functions}
\IMG{1103030806.pdf}
\LANGen{
\Question{
The function \( f \) is given by the graph. Identify which of the following statements is false.}
{The function \( f \) is non-increasing in the interval \( [ -3;2 ] \).}{The function \( f \) is not increasing.}{The function \( f \) is decreasing in the interval \( [ 2;5 ] \).}{The function \( f \) is non-decreasing in the interval  \( [ -1;2 ] \).}{}{}}
\LANGpl{
\Question{
Funkcja \( f \) przedstawiona jest za pomocą wykresu. Które z poniższych stwierdzeń jest fałszywe?}
{Funkcja \( f \) jest nierosnąca w przedziale \( \langle -3;2 \rangle \).}{Funkcja \( f \) nie jest rosnąca.}{Funkcja  \( f \) jest malejąca w przedziale \( \langle 2;5 \rangle \).}{Funkcja \( f \) jest niemalejąca w przedziale  \( \langle -1;2 \rangle \).}{}{}}
\LANGcs{
\Question{
Funkce \( f \) je dána grafem. Vyberte nepravdivý výrok.}
{Funkce \( f \) je v intervalu \( \langle -3;2 \rangle \) nerostoucí.}{Funkce \( f \) není rostoucí.}{Funkce \( f \) je v intervalu \( \langle 2;5 \rangle \) klesající.}{Funkce \( f \) je v intervalu  \( \langle -1;2 \rangle \) neklesající.}{}{}}
\LANGsk{
\Question{
Na obrázku je graf funkcie \( f \). Ktoré z nasledujúcich tvrdení o funkcii \( f \) je nepravdivé?}
{Funkcia \( f \) je na intervale \( \langle -3;2 \rangle \) nerastúca.}{Funkcia \( f \) nie je rastúca.}{Funkcia \( f \) je na intervale \( \langle 2;5 \rangle \) klesajúca.}{Funkcia \( f \) je na intervale \( \langle -1;2 \rangle \) neklesajúca.}{}{}}
\LANGes{
\Question{
La función \( f \) viene dada por la gráfica. ¿Cuál de las declaraciones es falsa?}
{La función \( f \) es no-creciente en el intervalo \( [ -3;2 ] \).}{La función \( f \) no es creciente.}{La función \( f \) es decreciente en el intervalo \( [ 2;5 ] \).}{La función \( f \) es no-decreciente en el intervalo \( [ -1;2 ] \).}{}{}}
\End

\Data\ID{1103030901}
\Updated{2020-12-30 10:12:46}
\Level{A}
\SubArea{Properties of functions}
\IMG{1103030901.pdf}
\LANGen{
\Question{
The function \( f \) is given by the following  graph. Identify which of the following statements is false.}
{The function \( f \) is not bounded below.}{The function \( f \) is bounded below.}{The function \( f \) is bounded above.}{The function \( f \) is bounded.}{}{}}
\LANGpl{
\Question{
Funkcję \( f \) przedstawiono za pomocą wykresu. Które z poniższych zdań jest fałszywe?}
{Funkcja \( f \) nie jest ograniczona z dołu.}{Funkcja \( f \) jest ograniczona z dołu.}{Funkcja  \( f \) jest ograniczona z góry.}{Funkcja \( f \) jest ograniczona.}{}{}}
\LANGcs{
\Question{
Funkce \( f \) je dána grafem. Vyberte nepravdivý výrok.}
{Funkce \( f \) není zdola omezená.}{Funkce  \( f \) je zdola omezená.}{Funkce \( f \) je shora omezená.}{Funkce \( f \) je omezená.}{}{}}
\LANGsk{
\Question{
Funkcia \( f \) je daná nasledujúcim grafom. Vyberte  nepravdivé tvrdenie o funkcii \( f \).}
{Funkcia \( f \) nie je zdola ohraničená.}{Funkcia \( f \) je zdola ohraničená.}{Funkcia \( f \) je zhora ohraničená.}{Funkcia \( f \) je ohraničená.}{}{}}
\LANGes{
\Question{
La función \( f \) viene dada por la gráfica. Idetifica cuál de las declaraciones siguientes es falsa.}
{La función \( f \) no es cerrada inferiormente.}{La función \( f \) es cerrada inferiormente.}{La función \( f \) es cerrada superiormente.}{La función \( f \) es cerrada.}{}{}}
\End

\Data\ID{1103048503}
\Updated{2020-12-30 10:12:03}
\Level{A}
\SubArea{Properties of functions}
\IMG{1103048503.pdf}
\LANGen{
\Question{
Let  \( f \) be a periodic function with period \( 4 \). The diagram shows a part of the graph of \( f \). Identify which of the following statements is false.}
{The function \( f \) is an odd function.}{The function \( f \) is increasing in the interval \( [14;15] \).}{The function \( f \) has maximum at \( x=-5 \).}{The function \( f \) is a bounded function.}{}{}}
\LANGpl{
\Question{
Załóżmy, że funkcja  \( f \) jest funkcją okresową z okresem \( 4 \). Poniższy diagram przedstawia część wykresu funkcji \( f \). Które z poniższych zdań jest  fałszywe?}
{Funkcja \( f \) jest funkcją nieparzystą.}{Funkcja \( f \) jest rosnąca w przedziale \( \langle14;15\rangle \).}{Funkcja \( f \) osiąga maksimum dla \( x=-5 \).}{Funkcja \( f \) jest ograniczona.}{}{}}
\LANGcs{
\Question{
Funkce \( f \)  je periodická s periodou  \( 4 \). Na obrázku je část jejího grafu. Vyberte nepravdivý výrok.}
{Funkce \( f \) je lichá.}{Funkce \( f \) je rostoucí v intervalu  \( \langle14;15\rangle \).}{Funkce \( f \) má maximum v bodě  \( x=-5 \).}{Funkce \( f \) je omezená.}{}{}}
\LANGsk{
\Question{
Daná je periodická funkcia \( f \) s periódou \( 4 \). Na obrázku je časť grafu funkcie \( f \). Vyberte nepravdivý výrok.}
{Funkcia \( f \) je nepárna.}{Funkcia \( f \) je rastúca na intervale \( \langle14;15\rangle \).}{Funkcia \( f \) má maximum v bode \( x=-5 \).}{Funkcia \( f \) je ohraničená.}{}{}}
\LANGes{
\Question{
Sea  \( f \) una función periódica con  periodo \( 4 \). En el diagrama se ve parte de la gráfica de \( f \). Identifica cuál de las declaraciones es falsa.}
{La función \( f \) es impar.}{La función \( f \) es creciente en el intervalo \( [14;15] \).}{La función \( f \) tiene máximo en \( x=-5 \).}{La función \( f \) es acotada.}{}{}}
\End

\Data\ID{9100004803}
\Updated{2020-09-02 03:09:11}
\Level{A}
\SubArea{Properties of functions}
\LANGen{
\Question{
Identify a possible graph of an even function.}
{}{}{}{}{}{}}
\LANGpl{
\Question{
Który obrazek przedstawia wykres funkcji parzystej?}
{}{}{}{}{}{}}
\LANGcs{
\Question{
Na kterém obrázku je znázorněn graf sudé funkce?}
{}{}{}{}{}{}}
\LANGsk{
\Question{
Na ktorom obrázku je znázornený graf párnej funkcie?}
{}{}{}{}{}{}}
\LANGes{
\Question{
NA}
{}{}{}{}{}{}}

\IMGanswer{000048_03-figure2.pdf}
\IMGanswer{000048_03-figure0.pdf}
\IMGanswer{000048_03-figure1.pdf}
\IMGanswer{000048_03-figure3.pdf}\End

\Data\ID{9100004806}
\Updated{2020-09-02 03:09:47}
\Level{A}
\SubArea{Properties of functions}
\LANGen{
\Question{
Identify a possible graph of an odd function.}
{}{}{}{}{}{}}
\LANGpl{
\Question{
Który obrazek przedstawia możliwy wykres funkcji nieparzystej?}
{}{}{}{}{}{}}
\LANGcs{
\Question{
Na kterém obrázku je znázorněn graf liché funkce?}
{}{}{}{}{}{}}
\LANGsk{
\Question{
Na ktorom obrázku je znázornený graf nepárnej funkcie?}
{}{}{}{}{}{}}
\LANGes{
\Question{
NA}
{}{}{}{}{}{}}

\IMGanswer{000048_06-figure3.pdf}
\IMGanswer{000048_06-figure0.pdf}
\IMGanswer{000048_06-figure1.pdf}
\IMGanswer{000048_06-figure2.pdf}\End

\Data\ID{1003019301}
\Updated{2020-09-01 11:09:44}
\Level{B}
\SubArea{Properties of functions}
\LANGen{
\Question{
Which of the following functions has maximum at \( x=-3 \)?}
{\( f(x) = -2x-3;\ x\in[-3;3)\)}{\( h(x) = 2x+3;\ x\in[-3;3]\)}{\( g(x) = -3x+2;\ x\in(-3;3]\)}{\( m(x) = 2x-3;\ x\in(-\infty;0]\)}{}{}}
\LANGpl{
\Question{
Która z poniższych funkcji ma maksimum w \( x=-3 \)?}
{\( f(x) = -2x-3;\ x\in\langle-3;3)\)}{\( h(x) = 2x+3;\ x\in\langle-3;3\rangle\)}{\( g(x) = -3x+2;\ x\in(-3;3\rangle\)}{\( m(x) = 2x-3;\ x\in(-\infty;0\rangle\)}{}{}}
\LANGcs{
\Question{
Která z následujících funkcí má maximum v bodě -3?}
{\( f(x) = -2x-3;\ x\in\langle-3;3)\)}{\( h(x) = 2x+3;\ x\in\langle-3;3\rangle\)}{\( g(x) = -3x+2;\ x\in(-3;3\rangle\)}{\( m(x) = 2x-3;\ x\in(-\infty;0\rangle\)}{}{}}
\LANGsk{
\Question{
Ktorá z následujúcich funkcií má maximum v bode -3?}
{\( f(x) = -2x-3;\ x\in\langle-3;3)\)}{\( h(x) = 2x+3;\ x\in\langle-3;3\rangle\)}{\( g(x) = -3x+2;\ x\in(-3;3\rangle\)}{\( m(x) = 2x-3;\ x\in(-\infty;0\rangle\)}{}{}}
\LANGes{
\Question{
¿Cuál de las funciones siguientes tiene máximo en \( x=-3 \)?}
{\( f(x) = -2x-3;\ x\in[-3;3)\)}{\( h(x) = 2x+3;\ x\in[-3;3]\)}{\( g(x) = -3x+2;\ x\in(-3;3]\)}{\( m(x) = 2x-3;\ x\in(-\infty;0]\)}{}{}}
\End

\Data\ID{1003028403}
\Updated{2020-09-01 12:09:35}
\Level{B}
\SubArea{Properties of functions}
\LANGen{
\Question{
Which of the following functions has the range \( [-3;\infty) \)?}
{\( f(x)= x^2-3\)}{\( h(x)= \sqrt{x-3} \)}{\( g(x)= |x+3| \)}{\( m(x)=\frac1x-3 \)}{}{}}
\LANGpl{
\Question{
Która z podanych funkcji ma zakres \( \langle-3;\infty) \)?}
{\( f(x)= x^2-3\)}{\( h(x)= \sqrt{x-3} \)}{\( g(x)= |x+3| \)}{\( m(x)=\frac1x-3 \)}{}{}}
\LANGcs{
\Question{
Vyberte funkci, jejímž oborem hodnot je interval \( \langle-3;\infty) \).}
{\( f(x)= x^2-3\)}{\( h(x)= \sqrt{x-3} \)}{\( g(x)= |x+3| \)}{\( m(x)=\frac1x-3 \)}{}{}}
\LANGsk{
\Question{
Ktorá z nasledujúcich funkcií má obor hodnôt určený intervalom \( \langle-3;\infty) \)?}
{\( f(x)= x^2-3\)}{\( h(x)= \sqrt{x-3} \)}{\( g(x)= |x+3| \)}{\( m(x)=\frac1x-3 \)}{}{}}
\LANGes{
\Question{
¿Cuál de las funciones siguientes tiene el Rango \( [-3;\infty) \)?}
{\( f(x)= x^2-3\)}{\( h(x)= \sqrt{x-3} \)}{\( g(x)= |x+3| \)}{\( m(x)=\frac1x-3 \)}{}{}}
\End

\Data\ID{1003028404}
\Updated{2020-12-30 10:12:25}
\Level{B}
\SubArea{Properties of functions}
\LANGen{
\Question{
Let \( f(x)=\frac{\sqrt{x+3}}{x^2-25} \). Which of the statements about the domain of the function \( f \) is true?}
{\( D(f)=[-3; 5)\cup (5;\infty) \)}{\( D(f)=(-3;5)\cup(5;\infty) \)}{\( D(f)=(-\infty;-5)\cup(-5;5)\cup(5;\infty) \)}{\( D(f)=(-\infty;-5)\cup(-5;-3)\cup(-3;5)\cup(5;\infty) \)}{}{}}
\LANGpl{
\Question{
Niech \( f(x)=\frac{\sqrt{x+3}}{x^2-25} \). Które ze stwierdzeń dotyczących dziedziny funkcji  \( f \) jest prawdziwe?}
{\( D(f)=\langle-3; 5)\cup (5;\infty) \)}{\( D(f)=(-3;5)\cup(5;\infty) \)}{\( D(f)=(-\infty;-5)\cup(-5;5)\cup(5;\infty) \)}{\( D(f)=(-\infty;-5)\cup(-5;-3)\cup(-3;5)\cup(5;\infty) \)}{}{}}
\LANGcs{
\Question{
Funkce f je dána předpisem \( f(x)=\frac{\sqrt{x+3}}{x^2-25} \). Vyberte pravdivý výrok o definičním oboru \(D(f)\) funkce f.}
{\( D(f)=\langle-3; 5)\cup (5;\infty) \)}{\( D(f)=(-3;5)\cup(5;\infty) \)}{\( D(f)=(-\infty;-5)\cup(-5;5)\cup(5;\infty) \)}{\( D(f)=(-\infty;-5)\cup(-5;-3)\cup(-3;5)\cup(5;\infty) \)}{}{}}
\LANGsk{
\Question{
Daná je funkcia  \( f(x)=\frac{\sqrt{x+3}}{x^2-25} \). Ktoré tvrdenie o definičnom obore funkcie \( f \) je pravdivé?}
{\( D(f)=\langle-3; 5)\cup (5;\infty) \)}{\( D(f)=(-3;5)\cup(5;\infty) \)}{\( D(f)=(-\infty;-5)\cup(-5;5)\cup(5;\infty) \)}{\( D(f)=(-\infty;-5)\cup(-5;-3)\cup(-3;5)\cup(5;\infty) \)}{}{}}
\LANGes{
\Question{
Sea \( f(x)=\frac{\sqrt{x+3}}{x^2-25} \). ¿Cuál de las declaraciones sobre el dominio de la función \( f \) es correcta?}
{\( D(f)=[-3; 5)\cup (5;\infty) \)}{\( D(f)=(-3;5)\cup(5;\infty) \)}{\( D(f)=(-\infty;-5)\cup(-5;5)\cup(5;\infty) \)}{\( D(f)=(-\infty;-5)\cup(-5;-3)\cup(-3;5)\cup(5;\infty) \)}{}{}}
\End

\Data\ID{1003028407}
\Updated{2020-12-30 10:12:03}
\Level{B}
\SubArea{Properties of functions}
\LANGen{
\Question{
Paul went by car from Ostrava to Olomouc for a business trip. There he spent \( 50 \) minutes at the meeting and then he went back the same way. Paul covered the distance of \( 98\,\mathrm{km} \) from Ostrava to Olomouc in \( 64 \) minutes. The distance back he covered in \( 66 \) minutes. Suppose the recording of the travelled distance and the time spent on the business trip started when Paul left Ostrava.  Dependence of this distance on the time describes the function \( s(t) \). Distance is in kilometres and time is in hours. Which of the following statements about the domain and the range of the function \( s \) is correct?}
{\( D(s)=[0;3] ;  H(s)=[0;196] \)}{\( D(s)=[0;196] ; H(s)=[0;3] \)}{\( D(s)=[0;3] ; H(s)=[0;98] \)}{\( D(s)=\left[0;\frac{13}6\right] ; H(s)=[0;196] \)}{}{}}
\LANGpl{
\Question{
Paweł pojechał samochodem z Ostrawy do Ołomuńca w podróż służbową. Tam spędził  \( 50 \) minut na spotkaniu, a następnie wrócił w ten sam sposób. Paweł pokonał dystans \( 98\,\mathrm{km} \)  z Ostrawy do Ołomuńca w \( 64 \) minuty. Odległość z powrotem pokonał w \( 66 \) minut. Załóżmy, że nagrywanie pokonanej odległości i czasu spędzonego w podróży służbowej rozpoczęło się, gdy Paweł opuścił Ostrawę. Zależność tej odległości od czasu opisuje funkcja \( s(t) \). Odległość podano w kilometrach, a czas w godzinach. Które z poniższych stwierdzeń dotyczących dziedziny i zakresu funkcji \( s \)  jest poprawne?}
{\( D(s)=\langle0;3\rangle ; H(s)=\langle0;196\rangle \)}{\( D(s)=\langle0;196\rangle ; H(s)=\langle0;3\rangle \)}{\( D(s)=\langle0;3\rangle ; H(s)=\langle0;98\rangle \)}{\( D(s)=\left\langle0;\frac{13}6\right\rangle ; H(s)=\langle0;196\rangle \)}{}{}}
\LANGcs{
\Question{
Pavel jel autem z Ostravy do Olomouce na služební cestu. Tam strávil 50 minut na jednání a poté se vydal stejnou cestou zpět. Cestu z Ostravy do Olomouce dlouhou 98 km ujel Pavel za 64 minut. Cesta zpátky mu trvala 66 minut. Předpokládejme, že dráhu auta a čas strávený na služební cestě začínáme měřit při výjezdu Pavla z Ostravy. Závislost uražené dráhy auta na čase stráveném na služební cestě popisuje funkce \( s(t)\). Hodnoty dráhy jsou udávány v kilometrech a hodnoty času v hodinách. Vyberte pravdivý výrok o definičním oboru \( D(s) \)  a oboru hodnot \( H(s) \) funkce \( s(t)\).}
{\( D(s)=\langle0;3\rangle ; H(s)=\langle0;196\rangle \)}{\( D(s)=\langle0;196\rangle ; H(s)=\langle0;3\rangle \)}{\( D(s)=\langle0;3\rangle ; H(s)=\langle0;98\rangle \)}{\( D(s)=\left\langle0;\frac{13}6\right\rangle ; H(s)=\langle0;196\rangle \)}{}{}}
\LANGsk{
\Question{
Peter cestoval autom z Ostravy do Olomouca na služobnú cestu. Strávil tam na stretnutí \(50\) minút a rovnakou cestou sa vrátil späť. Cesta z Ostravy do Olomouca dlhá \(98\,\mathrm{km}\) Petrovi trvala \(64\) minút. Cesta späť mu trvala \(66\) minút. Prepokladajme, že dráhu auta a čas strávený na služobnej ceste začíname merať pri odchode Petra z Ostravy. Závislosť prejdenej dráhy auta na čase strávenom na služobnej ceste popisuje funkcia \(s(t)\). Dráhu udávame v kilometroch a čas v hodinách. Ktoré tvrdenie o definičnom obore \(D(s)\) a obore hodnôt \(H(s)\) funkcie \(s(t)\) je pravdivé?}
{\( D(s)=\langle0;3\rangle ; H(s)=\langle0;196\rangle \)}{\( D(s)=\langle0;196\rangle ; H(s)=\langle0;3\rangle \)}{\( D(s)=\langle0;3\rangle ; H(s)=\langle0;98\rangle \)}{\( D(s)=\left\langle0;\frac{13}6\right\rangle ; H(s)=\langle0;196\rangle \)}{}{}}
\LANGes{
\Question{
Pablo fue de Ostrava a Olomouc en coche para un viaje de trabajo. La reunión de trabajo duró \( 50 \)  minutos y  volvió por la misma ruta. Pablo recorrió la distancia de \( 98\,\mathrm{km} \) de Ostrava a Olomouc en \( 64 \) minutos. Tardó \( 66 \) minutos en volver. Suponemos que tanto la distancia recorrida como el tiempo transcurrido viajando comienzan cuando Pablo salió de Ostrava. La función \( s(t) \) describe la dependencia de la distancia con respecto al tiempo. La distancia se mide en metros y el tiempo en horas. ¿Cuál de las declaraciones sobre el dominio y el rango de la función \( s \) es correcta?}
{\( D(s)=[0;3] ;  H(s)=[0;196] \)}{\( D(s)=[0;196] ; H(s)=[0;3] \)}{\( D(s)=[0;3] ; H(s)=[0;98] \)}{\( D(s)=\left[0;\frac{13}6\right] ; H(s)=[0;196] \)}{}{}}
\End

\Data\ID{1003030804}
\Updated{2020-09-01 02:09:03}
\Level{B}
\SubArea{Properties of functions}
\LANGen{
\Question{
Identify which of the following functions is decreasing in the interval \( [ 0;\infty ) \).}
{\( m(x)=-x^2 \)}{\( h(x) = \frac1x \)}{\( g(x) = -100+0.1x \)}{\( f(x) = (-x)^2 \)}{}{}}
\LANGpl{
\Question{
Która z podanych funkcji jest malejąca w przedziale \( \langle 0;\infty ) \)?}
{\( m(x)=-x^2 \)}{\( h(x) = \frac1x \)}{\( g(x) = -100+0{,}1x \)}{\( f(x) = (-x)^2 \)}{}{}}
\LANGcs{
\Question{
Která z následujících funkcí je klesající v intervalu \( \langle 0;\infty ) \)?}
{\( m(x)=-x^2 \)}{\( h(x) = \frac1x \)}{\( g(x) = -100+0{,}1x \)}{\( f(x) = (-x)^2 \)}{}{}}
\LANGsk{
\Question{
Ktorá z následujúcich funkcií je klesajúca na intervale \( \langle 0;\infty ) \)?}
{\( m(x)=-x^2 \)}{\( h(x) = \frac1x \)}{\( g(x) = -100+0{,}1x \)}{\( f(x) = (-x)^2 \)}{}{}}
\LANGes{
\Question{
¿Cuál de las funciones siguientes decrece en el intervalo \( [ 0;\infty ) \)?}
{\( m(x)=-x^2 \)}{\( h(x) = \frac1x \)}{\( g(x) = -100+0.1x \)}{\( f(x) = (-x)^2 \)}{}{}}
\End

\Data\ID{1003030805}
\Updated{2020-09-01 02:09:51}
\Level{B}
\SubArea{Properties of functions}
\LANGen{
\Question{
Identify which of the following functions is increasing in the interval  \( [ 1;\infty ) \).}
{\( m(x)=(1-x)^2 \)}{\( g(x)=1-x \)}{\( h(x)=1-x^2 \)}{\( f(x)=1-x^3 \)}{}{}}
\LANGpl{
\Question{
Która z poniższych funkcji jest rosnąca w przedziale \( \langle 1;\infty ) \)?}
{\( m(x)=(1-x)^2 \)}{\( g(x)=1-x \)}{\( h(x)=1-x^2 \)}{\( f(x)=1-x^3 \)}{}{}}
\LANGcs{
\Question{
Která z následujících funkcí je rostoucí v intervalu \( \langle 1;\infty ) \)?}
{\( m(x)=(1-x)^2 \)}{\( g(x)=1-x \)}{\( h(x)=1-x^2 \)}{\( f(x)=1-x^3 \)}{}{}}
\LANGsk{
\Question{
Ktorá z následujúcich funkcií je rastúca na intervale \( \langle 1;\infty ) \)?}
{\( m(x)=(1-x)^2 \)}{\( g(x)=1-x \)}{\( h(x)=1-x^2 \)}{\( f(x)=1-x^3 \)}{}{}}
\LANGes{
\Question{
¿Cuál de las funciones crece en el intervalo \( [ 1;\infty ) \)?}
{\( m(x)=(1-x)^2 \)}{\( g(x)=1-x \)}{\( h(x)=1-x^2 \)}{\( f(x)=1-x^3 \)}{}{}}
\End

\Data\ID{1003030807}
\Updated{2020-09-01 02:09:14}
\Level{B}
\SubArea{Properties of functions}
\LANGen{
\Question{
The function \( f(x) \) is increasing in the interval \( J \). Identify which of the following statements is false.}
{The function  \( h(x) = -2 f(x) \) is increasing in the interval \( J \).}{The function  \( g(x) = 2 f(x) \) is increasing in the interval \( J \).}{The function  \( m(x) = f(x)+2 \) is increasing in the interval \( J \).}{The function  \( n(x) = f(x)-2 \) is increasing in the interval \( J \).}{}{}}
\LANGpl{
\Question{
Funkcja \( f(x) \) jest rosnąca w przedziale \( J \). Które z podanych stwierdzeń jest fałszywe?}
{Funkcja   \( h(x) = -2 f(x) \) jest rosnąca w przedziale \( J \).}{Funkcja   \( g(x) = 2 f(x) \) jest rosnąca w przedziale \( J \).}{Funkcja   \( m(x) = f(x)+2 \)  jest rosnąca w przedziale \( J \).}{Funkcja   \( n(x) = f(x)-2 \)  jest rosnąca w przedziale \( J \).}{}{}}
\LANGcs{
\Question{
Funkce \( f(x) \) je rostoucí v intervalu \( J \). Vyberte nepravdivý výrok.}
{Funkce \( h(x) = -2 f(x) \) je rostoucí v intervalu \( J \).}{Funkce  \( g(x) = 2 f(x) \) je rostoucí v intervalu \( J \).}{Funkce  \( m(x) = f(x)+2 \) je rostoucí v intervalu \( J \).}{Funkce  \( n(x) = f(x)-2 \) je rostoucí v intervalu \( J \).}{}{}}
\LANGsk{
\Question{
Funkcia \( f(x) \) je rastúca na intervale \( J \). Ktoré z nasledujúcich tvrdení je nepravdivé?}
{Funkcia \( h(x) = -2 f(x) \) je rastúca na intervale \( J \).}{Funkcia  \( g(x) = 2 f(x) \) je rastúca na intervale \( J \).}{Funkcia  \( m(x) = f(x)+2 \) je rastúca na intervale \( J \).}{Funkcia \( n(x) = f(x)-2 \) je rastúca na intervale \( J \).}{}{}}
\LANGes{
\Question{
La función \( f(x) \) es creciente en el intervalo \( J \). ¿Cuál de las declaraciones siguientes es falsa?}
{La función  \( h(x) = -2 f(x) \) es creciente en el intervalo \( J \).}{La función  \( g(x) = 2 f(x) \) es creciente en el intervalo \( J \).}{La función  \( m(x) = f(x)+2 \) es creciente en el intervalo \( J \).}{La función  \( n(x) = f(x)-2 \) es creciente en el intervalo \( J \).}{}{}}
\End

\Data\ID{1003030904}
\Updated{2020-09-08 08:09:20}
\Level{B}
\SubArea{Properties of functions}
\LANGen{
\Question{
Identify which of the following functions is not bounded above.}
{\( m(x)=3-x^3 \)}{\( h(x)=3 \)}{\( g(x)=3-x^2 \)}{\( f(x)= 3-\sqrt x \)}{}{}}
\LANGpl{
\Question{
Określ, która z poniższych funkcji nie jest ograniczona z góry.}
{\( m(x)=3-x^3 \)}{\( h(x)=3 \)}{\( g(x)=3-x^2 \)}{\( f(x)= 3-\sqrt x \)}{}{}}
\LANGcs{
\Question{
Která z následujících funkcí není shora omezená?}
{\( m(x)=3-x^3 \)}{\( h(x)=3 \)}{\( g(x)=3-x^2 \)}{\( f(x)= 3-\sqrt x \)}{}{}}
\LANGsk{
\Question{
Ktorá z nasledujúcich funkcií nie je zhora ohraničená?}
{\( m(x)=3-x^3 \)}{\( h(x)=3 \)}{\( g(x)=3-x^2 \)}{\( f(x)= 3-\sqrt x \)}{}{}}
\LANGes{
\Question{
Identifica cuál de las funciones no es acotada superiormente.}
{\( m(x)=3-x^3 \)}{\( h(x)=3 \)}{\( g(x)=3-x^2 \)}{\( f(x)= 3-\sqrt x \)}{}{}}
\End

\Data\ID{1003048504}
\Updated{2020-09-08 08:09:40}
\Level{B}
\SubArea{Properties of functions}
\LANGen{
\Question{
Identify which of the following functions is not periodic.}
{\( m(x)=(-2)^x,x\in\mathbb{Z} \)}{\( h(x)=(-1)^x,x\in\mathbb{Z} \)}{\( g(x)=-2,x\in\mathbb{R}	 \)}{\( f(x)= 1^x,x\in\mathbb{R} \)}{}{}}
\LANGpl{
\Question{
Określ, która z podanych funkcji nie jest okresowa.}
{\( m(x)=(-2)^x,x\in\mathbb{Z} \)}{\( h(x)=(-1)^x,x\in\mathbb{Z} \)}{\( g(x)=-2,x\in\mathbb{R}	 \)}{\( f(x)= 1^x,x\in\mathbb{R} \)}{}{}}
\LANGcs{
\Question{
Která z následujících funkcí není periodická?}
{\( m(x)=(-2)^x,x\in\mathbb{Z} \)}{\( h(x)=(-1)^x,x\in\mathbb{Z} \)}{\( g(x)=-2,x\in\mathbb{R}	 \)}{\( f(x)= 1^x,x\in\mathbb{R} \)}{}{}}
\LANGsk{
\Question{
Ktorá z nasledujúcich funkcií nie je periodická.}
{\( m(x)=(-2)^x,x\in\mathbb{Z} \)}{\( h(x)=(-1)^x,x\in\mathbb{Z} \)}{\( g(x)=-2,x\in\mathbb{R}	 \)}{\( f(x)= 1^x,x\in\mathbb{R} \)}{}{}}
\LANGes{
\Question{
¿Cuál de las funciones no es periódica?}
{\( m(x)=(-2)^x,x\in\mathbb{Z} \)}{\( h(x)=(-1)^x,x\in\mathbb{Z} \)}{\( g(x)=-2,x\in\mathbb{R}	 \)}{\( f(x)= 1^x,x\in\mathbb{R} \)}{}{}}
\End

\Data\ID{1003048505}
\Updated{2020-09-08 08:09:48}
\Level{B}
\SubArea{Properties of functions}
\LANGen{
\Question{
Every real number \( x \) can be written as  \( x=c+d \), where \( c \) is an integer and \( d\in[0; 1) \). Then \( c \) is called the integer part of \( x \) and is denoted by \( [x] \). Which of the following functions has the biggest primitive period?}
{\( g(x)=(-1)^{[x]}  \)}{\( f(x)=[2x]-2x \)}{\( m(x)=3[x]-3x \)}{\( h(x)=[x]-x \)}{}{}}
\LANGpl{
\Question{
Każda liczba rzeczywista \( x \) może być zapisana jako  \( x=c+d \), gdzie \( c \) jest liczbą całkowitą i  \( d\in\langle0; 1) \). Zatem \( c \) zwane jest częścią całkowitą \( x \) i jest oznaczone przez \( [x] \). Która z poniższych funkcji ma największy okres podstawowy(zasadniczy)?}
{\( g(x)=(-1)^{[x]}  \)}{\( f(x)=[2x]-2x \)}{\( m(x)=3[x]-3x \)}{\( h(x)=[x]-x \)}{}{}}
\LANGcs{
\Question{
Každé reálné číslo \( x \) lze zapsat ve tvaru   \( x=c+d \), kde \( c \) je celé číslo a \( d\in\langle0; 1) \). Číslo \( c \) se nazývá celá část čísla \( x \) a označujeme je \( [x] \). Která z následujících funkcí má největší nejmenší periodu?}
{\( g(x)=(-1)^{[x]}  \)}{\( f(x)=[2x]-2x \)}{\( m(x)=3[x]-3x \)}{\( h(x)=[x]-x \)}{}{}}
\LANGsk{
\Question{
Každé reálne číslo \( x \) sa dá zapísať v tvare \( x=c+d \), kde \( c \) je celé číslo a \( d\in\langle0; 1) \). Číslo \( c \) sa nazýva celá časť čísla \( x \) a označujeme ho \( [x] \). Ktorá z nasledujúcich funkcií má najväčšiu najmenšiu periódu?}
{\( g(x)=(-1)^{[x]}  \)}{\( f(x)=[2x]-2x \)}{\( m(x)=3[x]-3x \)}{\( h(x)=[x]-x \)}{}{}}
\LANGes{
\Question{
Cada número real \( x \) puede escribirse como  \( x=c+d \), donde \( c \) es un número entero y \( d\in[0; 1) \). Luego a \( c \) le llamamos la parte entera de \( x \) y se denota \( [x] \). ¿Cuál de las funciones siguientes tiene el mayor periodo primitivo?}
{\( g(x)=(-1)^{[x]}  \)}{\( f(x)=[2x]-2x \)}{\( m(x)=3[x]-3x \)}{\( h(x)=[x]-x \)}{}{}}
\End

\Data\ID{1003048507}
\Updated{2020-09-08 08:09:22}
\Level{B}
\SubArea{Properties of functions}
\LANGen{
\Question{
Let \( f \) be  any periodic function. Identify which of the following statements is true.}
{The function \( f \) is not an injective (one-to-one) function.}{\( D(f)=\mathbb{R} \).}{The function \( f \) is bounded.}{The function \( f \) is increasing.}{}{}}
\LANGpl{
\Question{
Załóżmy, że funkcja \( f \) jest funkcją okresową. Określ, które z poniższych stwierdzeń jest prawdziwe.}
{Funkcja \( f \) nie jest iniekcją.}{\( D(f)=\mathbb{R} \).}{Funkcja  \( f \) jest ograniczona.}{Funkcja \( f \) jest rosnąca.}{}{}}
\LANGcs{
\Question{
Funkce \( f \) je libovolná periodická funkce. Vyberte pravdivý výrok.}
{Funkce \( f \) není prostá.}{\( D(f)=\mathbb{R} \).}{Funkce \( f \) je omezená.}{Funkce \( f \) je rostoucí.}{}{}}
\LANGsk{
\Question{
Daná je ľubovoľná periodická funkcia \( f \). Vyberte pravdivé tvrdenie o funkcii \( f \).}
{Funkcia \( f \) nie je prostá.}{\( D(f)=\mathbb{R} \).}{Funkcia \( f \) je ohraničená.}{Funkcia \( f \) je rastúca.}{}{}}
\LANGes{
\Question{
Sea \( f \) una función periódica. ¿Cuál de las declaraciones es correcta?}
{La función \( f \) no es inyectiva.}{\( D(f)=\mathbb{R} \).}{La función \( f \) es acotada.}{La función \( f \) es creciente.}{}{}}
\End

\Data\ID{9000004802}
\Updated{2020-09-02 01:09:54}
\Level{B}
\SubArea{Properties of functions}
\LANGen{
\Question{
In the following list identify an even function.}
{\(f(x)= |x|\)}{\(f(x) = |x + 1|\)}{\(f(x) = x + 1\)}{\(f(x) = x\)}{}{}}
\LANGpl{
\Question{
Która z niżej podanych funkcji jest funkcją parzystą?}
{\(f(x) = |x|\)}{\(f(x) = |x + 1|\)}{\(f(x) = x + 1\)}{\(f(x) = x\)}{}{}}
\LANGcs{
\Question{
Která z následujících funkcí je sudá?}
{\(f(x)= |x|\)}{\(f(x)= |x + 1|\)}{\(f(x) = x + 1\)}{\(f(x) = x\)}{}{}}
\LANGsk{
\Question{
Z daných funkcií vyberte párnu funkciu.}
{\(f(x) = |x|\)}{\(f(x) = |x + 1|\)}{\(f(x) = x + 1\)}{\(f(x) = x\)}{}{}}
\LANGes{
\Question{
¿Cuál de las funciones es par?}
{\(f(x)= |x|\)}{\(f(x) = |x + 1|\)}{\(f(x) = x + 1\)}{\(f(x) = x\)}{}{}}
\End

\Data\ID{9000004804}
\Updated{2020-09-02 01:09:13}
\Level{B}
\SubArea{Properties of functions}
\LANGen{
\Question{
In the following list identify an odd function.}
{\(f(x) = x^{3}\)}{\(f(x) = |x^{3}|\)}{\(f(x) = x^{4}\)}{\(f(x) = |x^{4}|\)}{}{}}
\LANGpl{
\Question{
Która z niżej podanych funkcji jest funkcją nieparzystą?}
{\(f(x) = x^{3}\)}{\(f(x) = |x^{3}|\)}{\(f(x)= x^{4}\)}{\(f(x) = |x^{4}|\)}{}{}}
\LANGcs{
\Question{
Která z následujících funkcí je lichá?}
{\(f(x) = x^{3}\)}{\(f(x) = |x^{3}|\)}{\(f(x) = x^{4}\)}{\(f(x) = |x^{4}|\)}{}{}}
\LANGsk{
\Question{
Z daných funkcií vyber nepárnu funkciu.}
{\(f(x) = x^{3}\)}{\(f(x) = |x^{3}|\)}{\(f(x) = x^{4}\)}{\(f(x) = |x^{4}|\)}{}{}}
\LANGes{
\Question{
¿Cuál de las funciones es impar?}
{\(f(x) = x^{3}\)}{\(f(x) = |x^{3}|\)}{\(f(x) = x^{4}\)}{\(f(x) = |x^{4}|\)}{}{}}
\End

\Data\ID{9000004805}
\Updated{2020-09-02 01:09:52}
\Level{B}
\SubArea{Properties of functions}
\LANGen{
\Question{
In the following list identify a function which is not an odd function.}
{\(y = x + 3\)}{\(y = x^{5}\)}{\(y = \frac{3} 
{x}\)}{\(y = x\)}{}{}}
\LANGpl{
\Question{
Która z niżej podanych funkcji nie jest funkcją nieparzystą?}
{\(y = x + 3\)}{\(y = x^{5}\)}{\(y = \frac{3} 
{x}\)}{\(y = x\)}{}{}}
\LANGcs{
\Question{
Která z následujících funkcí daných předpisem není lichá?}
{\(y = x + 3\)}{\(y = x^{5}\)}{\(y = \frac{3} 
{x}\)}{\(y = x\)}{}{}}
\LANGsk{
\Question{
Ktorá z nasledujúcich funkcií danej predpisom nie je nepárna?}
{\(y = x + 3\)}{\(y = x^{5}\)}{\(y = \frac{3} 
{x}\)}{\(y = x\)}{}{}}
\LANGes{
\Question{
¿Cuál de las funciones no es impar?}
{\(y = x + 3\)}{\(y = x^{5}\)}{\(y = \frac{3} 
{x}\)}{\(y = x\)}{}{}}
\End

\Data\ID{9000004809}
\Updated{2020-09-10 06:09:14}
\Level{B}
\SubArea{Properties of functions}
\LANGen{
\Question{
In the following list identify a function which is bounded above.}
{\(f(x) = -(x - 2)^{2}\)}{\(f(x) = x^{2}\)}{\(f(x) = x^{2} - 3x\)}{\(f(x) = (x - 3)^{2}\)}{}{}}
\LANGpl{
\Question{
Która z podanych funkcji jest ograniczona z góry?}
{\(f(x) = -(x - 2)^{2}\)}{\(f(x) = x^{2}\)}{\(f(x) = x^{2} - 3x\)}{\(f(x) = (x - 3)^{2}\)}{}{}}
\LANGcs{
\Question{
Která z následujících funkcí je omezená shora?}
{\(f(x) = -(x - 2)^{2}\)}{\(f(x) = x^{2}\)}{\(f(x) = x^{2} - 3x\)}{\(f(x) = (x - 3)^{2}\)}{}{}}
\LANGsk{
\Question{
Z nasledujúcich funkcií vyberte zhora ohraničenú funkciu.}
{\(f(x) = -(x - 2)^{2}\)}{\(f(x) = x^{2}\)}{\(f(x) = x^{2} - 3x\)}{\(f(x) = (x - 3)^{2}\)}{}{}}
\LANGes{
\Question{
¿Cuál de las funciones es acotada superiormente?}
{\(f(x) = -(x - 2)^{2}\)}{\(f(x) = x^{2}\)}{\(f(x) = x^{2} - 3x\)}{\(f(x) = (x - 3)^{2}\)}{}{}}
\End

\Data\ID{9000010601}
\Updated{2020-12-30 10:12:18}
\Level{B}
\SubArea{Properties of functions}
\LANGen{
\Question{
Identify a function which has a domain
\([ - 3;1] \).}
{\(y = \sqrt{-x^{2 } - 2x + 3}\)}{\(y = \sqrt{-x^{2 } + 2x - 3}\)}{\(y = \sqrt{x^{2 } + 2x - 3}\)}{\(y = \sqrt{x^{2 } - 2x + 3}\)}{\(y = \sqrt{\frac{x+3}
{x+1}}\)}{\(y = \sqrt{\frac{x-1}
{x+3}}\)}}
\LANGpl{
\Question{
Wybierz funkcję, której dziedziną jest przedział
\([ - 3;1] \).}
{\(y = \sqrt{-x^{2 } - 2x + 3}\)}{\(y = \sqrt{-x^{2 } + 2x - 3}\)}{\(y = \sqrt{x^{2 } + 2x - 3}\)}{\(y = \sqrt{x^{2 } - 2x + 3}\)}{\(y = \sqrt{\frac{x+3}
{x+1}}\)}{\(y = \sqrt{\frac{x-1}
{x+3}}\)}}
\LANGcs{
\Question{
Vyberte funkci, jejímž definičním oborem je interval
\(\langle - 3;1\rangle \).}
{\(y = \sqrt{-x^{2 } - 2x + 3}\)}{\(y = \sqrt{-x^{2 } + 2x - 3}\)}{\(y = \sqrt{x^{2 } + 2x - 3}\)}{\(y = \sqrt{x^{2 } - 2x + 3}\)}{\(y = \sqrt{\frac{x+3}
{x+1}}\)}{\(y = \sqrt{\frac{x-1}
{x+3}}\)}}
\LANGsk{
\Question{
Určte funkciu, ktorej definičný obor je interval
\(\langle - 3;1\rangle \).}
{\(y = \sqrt{-x^{2 } - 2x + 3}\)}{\(y = \sqrt{-x^{2 } + 2x - 3}\)}{\(y = \sqrt{x^{2 } + 2x - 3}\)}{\(y = \sqrt{x^{2 } - 2x + 3}\)}{\(y = \sqrt{\frac{x+3}
{x+1}}\)}{\(y = \sqrt{\frac{x-1}
{x+3}}\)}}
\LANGes{
\Question{
Elije la función cuyo dominio es
\([ - 3;1] \).}
{\(y = \sqrt{-x^{2 } - 2x + 3}\)}{\(y = \sqrt{-x^{2 } + 2x - 3}\)}{\(y = \sqrt{x^{2 } + 2x - 3}\)}{\(y = \sqrt{x^{2 } - 2x + 3}\)}{\(y = \sqrt{\frac{x+3}
{x+1}}\)}{\(y = \sqrt{\frac{x-1}
{x+3}}\)}}
\End

\Data\ID{9000010602}
\Updated{2020-12-30 10:12:48}
\Level{B}
\SubArea{Properties of functions}
\LANGen{
\Question{
Identify a function which has a domain
\((-\infty ;-2)\cup (2;\infty )\).}
{\(y = \sqrt{ \frac{1}
{x^{2}-4}}\)}{\(y = \frac{1} 
{x^{2}-4}\)}{\(y = \sqrt{x^{2 } + 4}\)}{\(y = \sqrt{x^{2 } - 2}\)}{\(y = \sqrt{x^{2 } - 4}\)}{\(y = \sqrt{ \frac{1}
{x^{2}-2}}\)}}
\LANGpl{
\Question{
Dziedziną, której z poniższych funkcji jest
\((-\infty ;-2)\cup (2;\infty )\)?}
{\(y = \sqrt{ \frac{1}
{x^{2}-4}}\)}{\(y = \frac{1} 
{x^{2}-4}\)}{\(y = \sqrt{x^{2 } + 4}\)}{\(y = \sqrt{x^{2 } - 2}\)}{\(y = \sqrt{x^{2 } - 4}\)}{\(y = \sqrt{ \frac{1}
{x^{2}-2}}\)}}
\LANGcs{
\Question{
Vyberte funkci, jejímž definičním oborem je množina
\(A = (-\infty ;-2)\cup (2;\infty )\).}
{\(y = \sqrt{ \frac{1}
{x^{2}-4}}\)}{\(y = \frac{1} 
{x^{2}-4}\)}{\(y = \sqrt{x^{2 } + 4}\)}{\(y = \sqrt{x^{2 } - 2}\)}{\(y = \sqrt{x^{2 } - 4}\)}{\(y = \sqrt{ \frac{1}
{x^{2}-2}}\)}}
\LANGsk{
\Question{
Určte funkciu, ktorej definičný obor je množina
\(A = (-\infty ;-2)\cup (2;\infty )\).}
{\(y = \sqrt{ \frac{1}
{x^{2}-4}}\)}{\(y = \frac{1} 
{x^{2}-4}\)}{\(y = \sqrt{x^{2 } + 4}\)}{\(y = \sqrt{x^{2 } - 2}\)}{\(y = \sqrt{x^{2 } - 4}\)}{\(y = \sqrt{ \frac{1}
{x^{2}-2}}\)}}
\LANGes{
\Question{
Elije la función cuyo dominio es
\((-\infty ;-2)\cup (2;\infty )\).}
{\(y = \sqrt{ \frac{1}
{x^{2}-4}}\)}{\(y = \frac{1} 
{x^{2}-4}\)}{\(y = \sqrt{x^{2 } + 4}\)}{\(y = \sqrt{x^{2 } - 2}\)}{\(y = \sqrt{x^{2 } - 4}\)}{\(y = \sqrt{ \frac{1}
{x^{2}-2}}\)}}
\End

\Data\ID{9000010603}
\Updated{2020-12-30 10:12:05}
\Level{B}
\SubArea{Properties of functions}
\LANGen{
\Question{
Identify a function which has a domain
\(\left (-\infty ;-\frac{3}
{2}\right ] \).}
{\(y = \sqrt{-2x - 3}\)}{\(y = \sqrt{3x + 2}\)}{\(y = -\sqrt{2 - 3x}\)}{\(y = \sqrt{x + \frac{3} 
{2}}\)}{\(y = \sqrt{x^{2 } - 3x}\)}{\(y = \frac{1} 
{3x+2}\)}}
\LANGpl{
\Question{
Dziedziną, której z poniższych funkcji jest
\(\left (-\infty ;-\frac{3}
{2}\right ] \)?}
{\(y = \sqrt{-2x - 3}\)}{\(y = \sqrt{3x + 2}\)}{\(y = -\sqrt{2 - 3x}\)}{\(y = \sqrt{x + \frac{3} 
{2}}\)}{\(y = \sqrt{x^{2 } - 3x}\)}{\(y = \frac{1} 
{3x+2}\)}}
\LANGcs{
\Question{
Vyberte funkci, jejímž definičním oborem je interval
\(\left (-\infty ;-\frac{3}
{2}\right \rangle \).}
{\(y = \sqrt{-2x - 3}\)}{\(y = \sqrt{3x + 2}\)}{\(y = -\sqrt{2 - 3x}\)}{\(y = \sqrt{x + \frac{3} 
{2}}\)}{\(y = \sqrt{x^{2 } - 3x}\)}{\(y = \frac{1} 
{3x+2}\)}}
\LANGsk{
\Question{
Určte funkciu, ktorej definičný obor je interval
\(\left (-\infty ;-\frac{3}
{2}\right \rangle \).}
{\(y = \sqrt{-2x - 3}\)}{\(y = \sqrt{3x + 2}\)}{\(y = -\sqrt{2 - 3x}\)}{\(y = \sqrt{x + \frac{3} 
{2}}\)}{\(y = \sqrt{x^{2 } - 3x}\)}{\(y = \frac{1} 
{3x+2}\)}}
\LANGes{
\Question{
Elije la función cuyo dominio es
\(\left (-\infty ;-\frac{3}
{2}\right ] \).}
{\(y = \sqrt{-2x - 3}\)}{\(y = \sqrt{3x + 2}\)}{\(y = -\sqrt{2 - 3x}\)}{\(y = \sqrt{x + \frac{3} 
{2}}\)}{\(y = \sqrt{x^{2 } - 3x}\)}{\(y = \frac{1} 
{3x+2}\)}}
\End

\Data\ID{9000010604}
\Updated{2020-12-30 10:12:19}
\Level{B}
\SubArea{Properties of functions}
\LANGen{
\Question{
Identify a function which has a domain
\([ - 3;5)\).}
{\(y = \sqrt{\frac{x+3}
{5-x}}\)}{\(y = \sqrt{(x - 3)(x + 5)}\)}{\(y = \sqrt{\frac{x-5}
{x+3}}\)}{\(y = \sqrt{(x - 5)(x + 3)}\)}{\(y =\log \frac{x+5} 
{x-3}\)}{\(y =\log \frac{x+3} 
{x-5}\)}}
\LANGpl{
\Question{
Wybierz funkcję, której dziedziną jest
\([ - 3;5)\).}
{\(y = \sqrt{\frac{x+3}
{5-x}}\)}{\(y = \sqrt{(x - 3)(x + 5)}\)}{\(y = \sqrt{\frac{x-5}
{x+3}}\)}{\(y = \sqrt{(x - 5)(x + 3)}\)}{\(y =\log \frac{x+5} 
{x-3}\)}{\(y =\log \frac{x+3} 
{x-5}\)}}
\LANGcs{
\Question{
Vyberte funkci, jejímž definičním oborem je interval
\(\langle - 3;5)\).}
{\(y = \sqrt{\frac{x+3}
{5-x}}\)}{\(y = \sqrt{(x - 3)(x + 5)}\)}{\(y = \sqrt{\frac{x-5}
{x+3}}\)}{\(y = \sqrt{(x - 5)(x + 3)}\)}{\(y =\log \frac{x+5} 
{x-3}\)}{\(y =\log \frac{x+3} 
{x-5}\)}}
\LANGsk{
\Question{
Určte funkciu, ktorej definičný obor je interval
\(\langle - 3;5)\).}
{\(y = \sqrt{\frac{x+3}
{5-x}}\)}{\(y = \sqrt{(x - 3)(x + 5)}\)}{\(y = \sqrt{\frac{x-5}
{x+3}}\)}{\(y = \sqrt{(x - 5)(x + 3)}\)}{\(y =\log \frac{x+5} 
{x-3}\)}{\(y =\log \frac{x+3} 
{x-5}\)}}
\LANGes{
\Question{
Elge la función que tiene el dominio
\([ - 3;5)\).}
{\(y = \sqrt{\frac{x+3}
{5-x}}\)}{\(y = \sqrt{(x - 3)(x + 5)}\)}{\(y = \sqrt{\frac{x-5}
{x+3}}\)}{\(y = \sqrt{(x - 5)(x + 3)}\)}{\(y =\log \frac{x+5} 
{x-3}\)}{\(y =\log \frac{x+3} 
{x-5}\)}}
\End

\Data\ID{9000021808}
\Updated{2020-12-30 10:12:23}
\Level{B}
\SubArea{Properties of functions}
\LANGen{
\Question{
Find the domain of the following function. 
\[ 
 f(x) = \sqrt{\frac{(x - 3)(x + 2)}
{(1 - x)(3 - x)}}
\]}
{\(\mathop{\mathrm{Dom}}(f) = (-\infty ;-2] \cup (1;3)\cup (3;\infty )\)}{\(\mathop{\mathrm{Dom}}(f) = (-\infty ;-2)\cup (1;3)\)}{\(\mathop{\mathrm{Dom}}(f) = (-\infty ;-2] \cup (1;\infty )\)}{\(\mathop{\mathrm{Dom}}(f) =[ -2;1)\cup (3;\infty )\)}{}{}}
\LANGpl{
\Question{
Znajdź dziedzinę podanej funkcji:
\[ 
 f\colon y = \sqrt{\frac{(x - 3)(x + 2)}
{(1 - x)(3 - x)}}
\]}
{\(\mathop{\mathrm{Dom}}(f) = (-\infty ;-2] \cup (1;3)\cup (3;\infty )\)}{\(\mathop{\mathrm{Dom}}(f) = (-\infty ;-2)\cup (1;3)\)}{\(\mathop{\mathrm{Dom}}(f) = (-\infty ;-2] \cup (1;\infty )\)}{\(\mathop{\mathrm{Dom}}(f) =[ -2;1)\cup (3;\infty )\)}{}{}}
\LANGcs{
\Question{
Definiční obor funkce \(f\colon y = \sqrt{\frac{(x-3)(x+2)}
{(1-x)(3-x)}}\) 
je:}
{\((-\infty ;-2\rangle \cup (1;3)\cup (3;\infty )\)}{\((-\infty ;-2)\cup (1;3)\)}{\((-\infty ;-2\rangle \cup (1;\infty )\)}{\(\langle - 2;1)\cup (3;\infty )\)}{}{}}
\LANGsk{
\Question{
Určte definičný obor funkcie \(f\colon y = \sqrt{\frac{(x-3)(x+2)}
{(1-x)(3-x)}}\).}
{\((-\infty ;-2\rangle \cup (1;3)\cup (3;\infty )\)}{\((-\infty ;-2)\cup (1;3)\)}{\((-\infty ;-2\rangle \cup (1;\infty )\)}{\(\langle - 2;1)\cup (3;\infty )\)}{}{}}
\LANGes{
\Question{
Encuentra el dominio de la siguiente función.
\[ 
 f(x) = \sqrt{\frac{(x - 3)(x + 2)}
{(1 - x)(3 - x)}}
\]}
{\(\mathop{\mathrm{Dom}}(f) = (-\infty ;-2] \cup (1;3)\cup (3;\infty )\)}{\(\mathop{\mathrm{Dom}}(f) = (-\infty ;-2)\cup (1;3)\)}{\(\mathop{\mathrm{Dom}}(f) = (-\infty ;-2] \cup (1;\infty )\)}{\(\mathop{\mathrm{Dom}}(f) =[ -2;1)\cup (3;\infty )\)}{}{}}
\End

\Data\ID{9000033703}
\Updated{2020-12-30 10:12:35}
\Level{B}
\SubArea{Properties of functions}
\LANGen{
\Question{
Find the domain of the following function. 
\[ 
 f\colon y = \frac{x} 
{\sqrt{4x^{2 } - 9}}
\]}
{\(\left (-\infty ;-\frac{3}
{2}\right )\cup \left (\frac{3}
{2};\infty \right )\)}{\(\mathbb{R}\)}{\(\mathbb{R}\setminus \left \{-\frac{3}
{2}; \frac{3} 
{2}\right \}\)}{\(\left (-\frac{3}
{2}; \frac{3} 
{2}\right )\)}{\(\left [ -\frac{3}
{2}; \frac{3} 
{2}\right ] \)}{\(\left (-\infty ;-\frac{3}
{2}\right ] \cup \left [ \frac{3}
{2};\infty \right )\)}}
\LANGpl{
\Question{
Znajdź dziedzinę podanej funkcji:
\[ 
 f\colon y = \frac{x} 
{\sqrt{4x^{2 } - 9}}
\]}
{\(\left (-\infty ;-\frac{3}
{2}\right )\cup \left (\frac{3}
{2};\infty \right )\)}{\(\mathbb{R}\)}{\(\mathbb{R}\setminus \left \{-\frac{3}
{2}; \frac{3} 
{2}\right \}\)}{\(\left (-\frac{3}
{2}; \frac{3} 
{2}\right )\)}{\(\left [ -\frac{3}
{2}; \frac{3} 
{2}\right ] \)}{\(\left (-\infty ;-\frac{3}
{2}\right ] \cup \left [ \frac{3}
{2};\infty \right )\)}}
\LANGcs{
\Question{
Definičním oborem funkce \(f\colon y = \frac{x} 
{\sqrt{4x^{2 } - 9}}\) 
je množina:}
{\(\left (-\infty ;-\frac{3}
{2}\right )\cup \left (\frac{3}
{2};\infty \right )\)}{\(\mathbb{R}\)}{\(\mathbb{R}\setminus \left \{-\frac{3}
{2}; \frac{3} 
{2}\right \}\)}{\(\left (-\frac{3}
{2}; \frac{3} 
{2}\right )\)}{\(\left \langle -\frac{3}
{2}; \frac{3} 
{2}\right \rangle \)}{\(\left (-\infty ;-\frac{3}
{2}\right \rangle \cup \left \langle \frac{3}
{2};\infty \right )\)}}
\LANGsk{
\Question{
Určte definičný obor funkcie \(f\colon y = \frac{x} 
{\sqrt{4x^{2 } - 9}}\).}
{\(\left (-\infty ;-\frac{3}
{2}\right )\cup \left (\frac{3}
{2};\infty \right )\)}{\(\mathbb{R}\)}{\(\mathbb{R}\setminus \left \{-\frac{3}
{2}; \frac{3} 
{2}\right \}\)}{\(\left (-\frac{3}
{2}; \frac{3} 
{2}\right )\)}{\(\left \langle -\frac{3}
{2}; \frac{3} 
{2}\right \rangle \)}{\(\left (-\infty ;-\frac{3}
{2}\right \rangle \cup \left \langle \frac{3}
{2};\infty \right )\)}}
\LANGes{
\Question{
Encuentra el dominio de la siguiente función.
\[ 
 f\colon y = \frac{x} 
{\sqrt{4x^{2 } - 9}}
\]}
{\(\left (-\infty ;-\frac{3}
{2}\right )\cup \left (\frac{3}
{2};\infty \right )\)}{\(\mathbb{R}\)}{\(\mathbb{R}\setminus \left \{-\frac{3}
{2}; \frac{3} 
{2}\right \}\)}{\(\left (-\frac{3}
{2}; \frac{3} 
{2}\right )\)}{\(\left [ -\frac{3}
{2}; \frac{3} 
{2}\right ] \)}{\(\left (-\infty ;-\frac{3}
{2}\right ] \cup \left [ \frac{3}
{2};\infty \right )\)}}
\End

\Data\ID{1003030401}
\Updated{2020-12-30 10:12:31}
\Level{C}
\SubArea{Properties of functions}
\LANGen{
\Question{
Suppose function \( f \) is given completely by the next  table.
\[
\begin{array}{|c|c|c|c|c|c|c|c|} \hline x&-3&-2&-1&0&1&2&3 \\\hline f(x)&-1&0&1&2&3&4&5 \\\hline
\end{array}\]
Identify which of the following functions is the inverse of \( f \).}
{Function \( h \), which is given completely by the next table.
\(
\begin{array}{|c|c|c|c|c|c|c|c|} \hline x&-1&0&1&2&3&4&5 \\\hline h(x)&-3&-2&-1&0&1&2&3 \\\hline
\end{array}\)}{Function \( m \), which is given completely by the next table.
\(\begin{array}{|c|c|c|c|c|c|c|c|} \hline x&-3&-2&-1&0&1&2&3 \\\hline m(x)&5&4&3&2&1&0&-1 \\\hline
\end{array}\)}{Function \( g \), such that \( g(x)=x-2 \) for \( x\in[-1;5] \).}{Function \( n \), such that \( n(x)=x+2 \) for \( x\in[-3;3] \).}{}{}}
\LANGpl{
\Question{
Załóżmy, że funkcja \( f \) jest całkowicie wyrażona w tabeli poniżej.
\[
\begin{array}{|c|c|c|c|c|c|c|c|} \hline x&-3&-2&-1&0&1&2&3 \\\hline f(x)&-1&0&1&2&3&4&5 \\\hline
\end{array}\]
Która z podanych funkcji jest odwrotnością funkcji \( f \)?}
{Funkcja \( h \), wyrażona w pełni w tabeli.
\(
\begin{array}{|c|c|c|c|c|c|c|c|} \hline x&-1&0&1&2&3&4&5 \\\hline h(x)&-3&-2&-1&0&1&2&3 \\\hline
\end{array}\)}{Funkcja  \( m \) wyrażona w pełni w tabeli.
\(\begin{array}{|c|c|c|c|c|c|c|c|} \hline x&-3&-2&-1&0&1&2&3 \\\hline m(x)&5&4&3&2&1&0&-1 \\\hline
\end{array}\)}{Funkcja  \( g \), taka, że  \( g(x)=x-2 \) dla \( x\in\langle-1;5\rangle \).}{Funkcja \( n \), taka, że \( n(x)=x+2 \) dla \( x\in\langle-3;3\rangle \).}{}{}}
\LANGcs{
\Question{
Funkce \( f \) je dána tabulkou:
\[
\begin{array}{|c|c|c|c|c|c|c|c|} \hline x&-3&-2&-1&0&1&2&3 \\\hline f(x)&-1&0&1&2&3&4&5 \\\hline
\end{array}\]
Která z následujících funkcí je funkcí inverzní k funkci \( f \)?}
{Funkce \( h \), která je dána tabulkou:
\(
\begin{array}{|c|c|c|c|c|c|c|c|} \hline x&-1&0&1&2&3&4&5 \\\hline h(x)&-3&-2&-1&0&1&2&3 \\\hline
\end{array}
\)}{Funkce \( m \), která je dána tabulkou:
\(
\begin{array}{|c|c|c|c|c|c|c|c|} \hline x&-3&-2&-1&0&1&2&3 \\\hline m(x)&5&4&3&2&1&0&-1 \\\hline
\end{array}
\)}{Funkce \( g(x)=x-2; \  x\in\langle-1;5\rangle \).}{Funkce \( n(x)=x+2 ;\ x\in\langle-3;3\rangle \).}{}{}}
\LANGsk{
\Question{
Predpokladajme, že tabuľka určuje úplnú funkciu \( f \). 
\[
\begin{array}{|c|c|c|c|c|c|c|c|} \hline x&-3&-2&-1&0&1&2&3 \\\hline f(x)&-1&0&1&2&3&4&5 \\\hline
\end{array}\]
Ktorá z následujúcich funkcií je inverznou funkciou k funkcii \( f \)?}
{Funkcia \( h \), ktorá je daná tabuľkou:
\(
\begin{array}{|c|c|c|c|c|c|c|c|} \hline x&-1&0&1&2&3&4&5 \\\hline h(x)&-3&-2&-1&0&1&2&3 \\\hline
\end{array}
\)}{Funkcia \( m \), ktorá je daná tabuľkou:
\(
\begin{array}{|c|c|c|c|c|c|c|c|} \hline x&-3&-2&-1&0&1&2&3 \\\hline m(x)&5&4&3&2&1&0&-1 \\\hline
\end{array}
\)}{Funkcia \( g(x)=x-2; \  x\in\langle-1;5\rangle \).}{Funkcia \( n(x)=x+2 ;\ x\in\langle-3;3\rangle \).}{}{}}
\LANGes{
\Question{
Supongamos que la función \( f \) viene dada por siguiente tabla:
\[
\begin{array}{|c|c|c|c|c|c|c|c|} \hline x&-3&-2&-1&0&1&2&3 \\\hline f(x)&-1&0&1&2&3&4&5 \\\hline
\end{array}\]
¿Cuál de las funciones siguiente es la inversa de \( f \)?}
{Función \( h \), dada por la tabla:
\(
\begin{array}{|c|c|c|c|c|c|c|c|} \hline x&-1&0&1&2&3&4&5 \\\hline h(x)&-3&-2&-1&0&1&2&3 \\\hline
\end{array}\)}{Función \( m \), dada por la tabla:
\(\begin{array}{|c|c|c|c|c|c|c|c|} \hline x&-3&-2&-1&0&1&2&3 \\\hline m(x)&5&4&3&2&1&0&-1 \\\hline
\end{array}\)}{Función \( g \), tal que \( g(x)=x-2 \) para \( x\in[-1;5] \).}{Función \( n \), tal que \( n(x)=x+2 \) para \( x\in[-3;3] \).}{}{}}
\End

\Data\ID{1003030402}
\Updated{2020-12-30 10:12:55}
\Level{C}
\SubArea{Properties of functions}
\LANGen{
\Question{
Suppose function \( f \) is given completely by the next table.
\[
\begin{array}{|c|c|c|c|c|c|c|c|} \hline x&-3&-2&-1&0&1&2&3 \\\hline f(x)&-1&2&-3&1&-2&3&2 \\\hline
\end{array}\]
Identify which of the following statements is true.}
{The inverse of \( f \) does not exist.}{The inverse of  \( f \) is function \( h \), which is given completely by the next table. 
\(
\begin{array}{|c|c|c|c|c|c|c|c|} \hline x&-1&2&-3&1&-2&3&2 \\\hline h(x)&-3&-2&-1&0&1&2&3 \\\hline
\end{array}
\)}{The inverse of  \( f \) is function \( g \), which is given completely by the next table. \(
\begin{array}{|c|c|c|c|c|c|c|c|} \hline x&-3&-2&-1&0&1&2&3 \\\hline g(x)&1&-2&3&-1&2&-3&-2 \\\hline
\end{array}\)}{The inverse of  \( f \) is function \( m \), which is given completely by the next table. \(
\begin{array}{|c|c|c|c|c|c|c|c|} \hline x&3&2&1&0&-1&-2&-3 \\\hline m(x)&1&-2&3&-1&2&-3&-2 \\\hline
\end{array}\)}{}{}}
\LANGpl{
\Question{
Załóżmy, że funkcja  \( f \) jest wyrażona całkowicie w podanej tabeli. 
\[
\begin{array}{|c|c|c|c|c|c|c|c|} \hline x&-3&-2&-1&0&1&2&3 \\\hline f(x)&-1&2&-3&1&-2&3&2 \\\hline
\end{array}\]
Które ze stwierdzeń jest prawdziwe?}
{Odwrotność funkcji  \( f \) nie istnieje.}{Odwrotnością funkcji \( f \) jest funkcja  \( h \), podana całkowicie w tabeli.
\(
\begin{array}{|c|c|c|c|c|c|c|c|} \hline x&-1&2&-3&1&-2&3&2 \\\hline h(x)&-3&-2&-1&0&1&2&3 \\\hline
\end{array}
\)}{Odwrotnością funkcji  \( f \) jest funkcja  \( g \), podana całkowicie w poniższej tabeli. \(
\begin{array}{|c|c|c|c|c|c|c|c|} \hline x&-3&-2&-1&0&1&2&3 \\\hline g(x)&1&-2&3&-1&2&-3&-2 \\\hline
\end{array}\)}{Odwrotnością funkcji \( f \) jest funkcja  \( m \),  podana całkowicie w poniższej tabeli. \(
\begin{array}{|c|c|c|c|c|c|c|c|} \hline x&3&2&1&0&-1&-2&-3 \\\hline m(x)&1&-2&3&-1&2&-3&-2 \\\hline
\end{array}\)}{}{}}
\LANGcs{
\Question{
Funkce \( f \) je dána tabulkou:\[
\begin{array}{|c|c|c|c|c|c|c|c|} \hline x&-3&-2&-1&0&1&2&3 \\\hline f(x)&-1&2&-3&1&-2&3&2 \\\hline
\end{array}\]
Vyberte pravdivý výrok.}
{K funkci \( f \) neexistuje inverzní funkce.}{Inverzní funkcí k funkci  \( f \)  je funkce \( h \), která je dána tabulkou: 
\(
\begin{array}{|c|c|c|c|c|c|c|c|} \hline x&-1&2&-3&1&-2&3&2 \\\hline h(x)&-3&-2&-1&0&1&2&3 \\\hline
\end{array}\)}{Inverzní funkcí k funkci  \( f \)  je funkce \( g \), která je dána tabulkou: \(
\begin{array}{|c|c|c|c|c|c|c|c|} \hline x&-3&-2&-1&0&1&2&3 \\\hline g(x)&1&-2&3&-1&2&-3&-2 \\\hline
\end{array}\)}{Inverzní funkcí k funkci  \( f \)  je funkce \( m \), která je dána tabulkou: \(
\begin{array}{|c|c|c|c|c|c|c|c|} \hline x&3&2&1&0&-1&-2&-3 \\\hline m(x)&1&-2&3&-1&2&-3&-2 \\\hline
\end{array}\)}{}{}}
\LANGsk{
\Question{
Predpokladajme, že tabuľka určuje úplnú funkciu \( f \). 
\[
\begin{array}{|c|c|c|c|c|c|c|c|} \hline x&-3&-2&-1&0&1&2&3 \\\hline f(x)&-1&2&-3&1&-2&3&2 \\\hline
\end{array}\]
Ktoré z nasledujúcich tvrdení je pravdivé?}
{K funkcii \( f \) neexistuje inverzná funkcia.}{Inverzná funkcia k funkcii  \( f \)  je funkcia \( h \), ktorá je daná tabuľkou: 
\(
\begin{array}{|c|c|c|c|c|c|c|c|} \hline x&-1&2&-3&1&-2&3&2 \\\hline h(x)&-3&-2&-1&0&1&2&3 \\\hline
\end{array}\)}{Inverzná funkcia k funkcii  \( f \)  je funkcia \( g \), ktorá je daná tabuľkou:  \(
\begin{array}{|c|c|c|c|c|c|c|c|} \hline x&-3&-2&-1&0&1&2&3 \\\hline g(x)&1&-2&3&-1&2&-3&-2 \\\hline
\end{array}\)}{Inverzná funkcia k funkcii  \( f \)  je funkcia \( m \), ktorá je daná tabuľkou:  \(
\begin{array}{|c|c|c|c|c|c|c|c|} \hline x&3&2&1&0&-1&-2&-3 \\\hline m(x)&1&-2&3&-1&2&-3&-2 \\\hline
\end{array}\)}{}{}}
\LANGes{
\Question{
Supongamos que la función \( f \) viene dada por siguiente tabla:
\[
\begin{array}{|c|c|c|c|c|c|c|c|} \hline x&-3&-2&-1&0&1&2&3 \\\hline f(x)&-1&2&-3&1&-2&3&2 \\\hline
\end{array}\]
¿Cuál de las declaraciones es correcta?}
{La inversa de \( f \) no existe.}{La inversa de  \( f \) es la función \( h \), dada por esta tabla:
\(
\begin{array}{|c|c|c|c|c|c|c|c|} \hline x&-1&2&-3&1&-2&3&2 \\\hline h(x)&-3&-2&-1&0&1&2&3 \\\hline
\end{array}
\)}{La inversa de  \( f \) es la función  \( g \), dada por esta tabla:
 \(
\begin{array}{|c|c|c|c|c|c|c|c|} \hline x&-3&-2&-1&0&1&2&3 \\\hline g(x)&1&-2&3&-1&2&-3&-2 \\\hline
\end{array}\)}{La inversa de  \( f \) es la función  \( m \),  dada por esta tabla:
 \(
\begin{array}{|c|c|c|c|c|c|c|c|} \hline x&3&2&1&0&-1&-2&-3 \\\hline m(x)&1&-2&3&-1&2&-3&-2 \\\hline
\end{array}\)}{}{}}
\End

\Data\ID{1003030403}
\Updated{2020-09-01 02:09:03}
\Level{C}
\SubArea{Properties of functions}
\LANGen{
\Question{
Identify which of the following functions is not an injective (one-to-one) function.}
{\( m(x)=|x-2|;\ x\in[-2; 3 ) \)}{\(  h(x)=x^3+3;\ x\in[-2;3) \)}{\( g(x)=(x+2)^2;\ x\in[-2; 3) \)}{\( f(x)=  \frac{x+2}{x-3};\ x\in[-2; 3) \)}{}{}}
\LANGpl{
\Question{
Która z podanych funkcji nie jest funkcją iniektywną (jeden do jednego)?}
{\( m(x)=|x-2|;\ x\in\langle-2; 3 ) \)}{\(  h(x)=x^3+3;\ x\in\langle-2;3) \)}{\( g(x)=(x+2)^2;\ x\in\langle-2; 3) \)}{\( f(x)=  \frac{x+2}{x-3};\ x\in\langle-2; 3) \)}{}{}}
\LANGcs{
\Question{
Která z následujících funkcí není prostá?}
{\( m(x)=|x-2|;\ x\in\langle-2; 3 ) \)}{\(  h(x)=x^3+3;\ x\in\langle-2;3) \)}{\( g(x)=(x+2)^2;\ x\in\langle-2; 3) \)}{\( f(x)=  \frac{x+2}{x-3};\ x\in\langle-2; 3) \)}{}{}}
\LANGsk{
\Question{
Ktorá z následujúcich funkcií nie je prostá?}
{\( m(x)=|x-2|;\ x\in\langle-2; 3 ) \)}{\(  h(x)=x^3+3;\ x\in\langle-2;3) \)}{\( g(x)=(x+2)^2;\ x\in\langle-2; 3) \)}{\( f(x)=  \frac{x+2}{x-3};\ x\in\langle-2; 3) \)}{}{}}
\LANGes{
\Question{
¿Cuál de las funciones no es inyectiva?}
{\( m(x)=|x-2|;\ x\in[-2; 3 ) \)}{\(  h(x)=x^3+3;\ x\in[-2;3) \)}{\( g(x)=(x+2)^2;\ x\in[-2; 3) \)}{\( f(x)=  \frac{x+2}{x-3};\ x\in[-2; 3) \)}{}{}}
\End

\Data\ID{1003030404}
\Updated{2020-09-01 02:09:11}
\Level{C}
\SubArea{Properties of functions}
\LANGen{
\Question{
Identify which of the following points lies on the graph of the inverse of function \( f(x)=x^3+2 \).}
{\( [-6;-2] \)}{\( [6;2] \)}{\( [-2;-6] \)}{\( [-6;2] \)}{\( [-2;6] \)}{}}
\LANGpl{
\Question{
Który z poniższych punktów należy do wykresu odwrotności funkcji \( f(x)=x^3+2 \)?}
{\( [-6;-2] \)}{\( [6;2] \)}{\( [-2;-6] \)}{\( [-6;2] \)}{\( [-2;6] \)}{}}
\LANGcs{
\Question{
Který z následujících bodů leží na grafu funkce inverzní k funkci \( f(x)=x^3+2 \)?}
{\( [-6;-2] \)}{\( [6;2] \)}{\( [-2;-6] \)}{\( [-6;2] \)}{\( [-2;6] \)}{}}
\LANGsk{
\Question{
Ktorý z následujúcich bodov leží na grafe funkcie inverznej k funkcii \( f(x)=x^3+2 \)?}
{\( [-6;-2] \)}{\( [6;2] \)}{\( [-2;-6] \)}{\( [-6;2] \)}{\( [-2;6] \)}{}}
\LANGes{
\Question{
Identifica cuál de los puntos pertenece a la gráfica de la inversa de la función \( f(x)=x^3+2 \).}
{\( [-6;-2] \)}{\( [6;2] \)}{\( [-2;-6] \)}{\( [-6;2] \)}{\( [-2;6] \)}{}}
\End

\Data\ID{1003030405}
\Updated{2020-09-01 02:09:09}
\Level{C}
\SubArea{Properties of functions}
\LANGen{
\Question{
Identify which of the following functions is the inverse of function \( f(x)=-\frac12 x+3 \).}
{\( h(x)=-2x+6 \)}{\( g(x)=-2x+3 \)}{\( m(x)=-2x+\frac13 \)}{\( n(x)=\frac12 x-3 \)}{}{}}
\LANGpl{
\Question{
Która z podanych funkcji jest odwrotnością funkcji  \( f(x)=-\frac12 x+3 \)?}
{\( h(x)=-2x+6 \)}{\( g(x)=-2x+3 \)}{\( m(x)=-2x+\frac13 \)}{\( n(x)=\frac12 x-3 \)}{}{}}
\LANGcs{
\Question{
Která z následujících funkcí je funkcí inverzní k funkci \( f(x)=-\frac12 x+3 \)?}
{\( h(x)=-2x+6 \)}{\( g(x)=-2x+3 \)}{\( m(x)=-2x+\frac13 \)}{\( n(x)=\frac12 x-3 \)}{}{}}
\LANGsk{
\Question{
Ktorá z následujúcich funkcií je inverznou funkciou k funkcii \( f(x)=-\frac12 x+3 \)?}
{\( h(x)=-2x+6 \)}{\( g(x)=-2x+3 \)}{\( m(x)=-2x+\frac13 \)}{\( n(x)=\frac12 x-3 \)}{}{}}
\LANGes{
\Question{
¿Cuál de las funciones es la inversa de la función \( f(x)=-\frac12 x+3 \)?}
{\( h(x)=-2x+6 \)}{\( g(x)=-2x+3 \)}{\( m(x)=-2x+\frac13 \)}{\( n(x)=\frac12 x-3 \)}{}{}}
\End

\Data\ID{1003030406}
\Updated{2020-09-01 02:09:03}
\Level{C}
\SubArea{Properties of functions}
\LANGen{
\Question{
Identify which of the following functions is the inverse of function \( f(x)=\frac{2x+1}{x-3 }\).}
{\( h(x)=\frac{3x+1}{x-2} \)}{\( g(x)=\frac{x-3}{2x+1} \)}{\( m(x)=\frac{2x-1}{x+3} \)}{\(  n(x)=\frac{3x-1}{x+2} \)}{}{}}
\LANGpl{
\Question{
Która z podanych funkcji jest odwrotnością funkcji  \( f(x)=\frac{2x+1}{x-3 }\)?}
{\( h(x)=\frac{3x+1}{x-2} \)}{\( g(x)=\frac{x-3}{2x+1} \)}{\( m(x)=\frac{2x-1}{x+3} \)}{\(  n(x)=\frac{3x-1}{x+2} \)}{}{}}
\LANGcs{
\Question{
Která z následujících funkcí je inverzní k funkci \( f(x)=\frac{2x+1}{x-3 }\)?}
{\( h(x)=\frac{3x+1}{x-2} \)}{\( g(x)=\frac{x-3}{2x+1} \)}{\( m(x)=\frac{2x-1}{x+3} \)}{\(  n(x)=\frac{3x-1}{x+2} \)}{}{}}
\LANGsk{
\Question{
Ktorá z následujúcich funkcií je inverznou funkciou k funkcii \( f(x)=\frac{2x+1}{x-3 }\)?}
{\( h(x)=\frac{3x+1}{x-2} \)}{\( g(x)=\frac{x-3}{2x+1} \)}{\( m(x)=\frac{2x-1}{x+3} \)}{\(  n(x)=\frac{3x-1}{x+2} \)}{}{}}
\LANGes{
\Question{
¿Cuál de las funciones es la inversa de la función \( f(x)=\frac{2x+1}{x-3 }\)?}
{\( h(x)=\frac{3x+1}{x-2} \)}{\( g(x)=\frac{x-3}{2x+1} \)}{\( m(x)=\frac{2x-1}{x+3} \)}{\(  n(x)=\frac{3x-1}{x+2} \)}{}{}}
\End

\Data\ID{1003107101}
\Updated{2020-09-08 09:09:34}
\Level{C}
\SubArea{Properties of functions}
\LANGen{
\Question{
Complete the following true statement: The graph of a function \( f(x) \) is the reflection the graph of its inverse \( f^{-1}(x) \)}
{across the line \( y=x \).}{across the axis \( x \).}{across the axis \( y \).}{about the origin.}{}{}}
\LANGpl{
\Question{
Uzupełnij tak, by stwierdzenie było prawdziwe: Wykres funkcji \( f(x) \) jest odbiciem jej odwrotności \( f^{-1}(x) \)}
{przez linię \( y=x \).}{przez współrzędną \( x \).}{przez współrzędną \( y \).}{o pochodzeniu.}{}{}}
\LANGcs{
\Question{
Doplňte následující pravdivý výrok: Grafy funkce \( f(x) \) a funkce k ní inverzní jsou}
{osově souměrné podle přímky \( y=x \).}{osově souměrné podle osy \( x \).}{osově souměrné podle osy \( y \).}{středově souměrné podle počátku souřadnicového systému.}{}{}}
\LANGsk{
\Question{
Doplňte nasledujúce pravdivé tvrdenie: Grafy funkcie \( f(x) \) a funkcie k nej inverznej \( f^{-1}(x) \) sú}
{osovo súmerné podľa priamky \( y=x \).}{osovo súmerné podľa osi \( x \).}{osovo súmerné podľa osi \( y \).}{stredovo súmerné podľa počiatku súradného systému.}{}{}}
\LANGes{
\Question{
Completa la declaración para que sea verdadera:  La gráfica de la función \( f(x) \) es la reflexión de la gráfica de su función inversa \( f^{-1}(x) \)}
{respecto a la recta \( y=x \).}{respecto al eje \( x \).}{respecto al eje \( y \).}{respecto al origen.}{}{}}
\End

\Data\ID{1103030407}
\Updated{2020-12-30 10:12:10}
\Level{C}
\SubArea{Properties of functions}
\IMG{1103030407.pdf}
\LANGen{
\Question{
Suppose function \( f \) is given by the graph. Identify which of the following graphs is the graph of the inverse of \( f \).}
{}{}{}{}{}{}}
\LANGpl{
\Question{
Załóżmy, że funkcja \( f \) jest wyrażona za pomocą wykresu.  Który wykres przedstawia odwrotność funkcji \( f \)?}
{}{}{}{}{}{}}
\LANGcs{
\Question{
Funkce $f$ je dána grafem. Na kterém obrázku je graf inverzní funkce k funkci \( f \)?}
{}{}{}{}{}{}}
\LANGsk{
\Question{
Funkcia \( f \) je daná grafom. Na ktorom obrázku je graf inverznej funkcie k funkcii \( f \)?}
{}{}{}{}{}{}}
\LANGes{
\Question{
NA}
{}{}{}{}{}{}}

\IMGanswer{1103030407a.pdf}
\IMGanswer{1103030407b.pdf}
\IMGanswer{1103030407c.pdf}
\IMGanswer{1103030407d.pdf}\End

\Data\ID{1103030408}
\Updated{2020-09-08 01:09:16}
\Level{C}
\SubArea{Properties of functions}
\LANGen{
\Question{
In one of the following pictures is the graph of the function \( f \), which \textbf{is not} one-to-one function. Choose the picture.}
{}{}{}{}{}{}}
\LANGpl{
\Question{
Jeden z poniższych rysunków przedstawia wykres funkcji \( f \), gdzie \textbf{nie jest } funkcją jeden do jednego. Wybierz odpowiedni rysunek.}
{}{}{}{}{}{}}
\LANGcs{
\Question{
Na jednom z následujících obrázků je graf funkce \( f \), která \textbf{není} prostá. Vyberte tento obrázek.}
{}{}{}{}{}{}}
\LANGsk{
\Question{
Na ktorom z nasledujúcich obrázkov je graf funkcie \( f \), ktorá \textbf{nie je} prostá.}
{}{}{}{}{}{}}
\LANGes{
\Question{
NA}
{}{}{}{}{}{}}

\IMGanswer{1103030408a.pdf}
\IMGanswer{1103030408b_0.pdf}
\IMGanswer{1103030408c_0.pdf}
\IMGanswer{1103030408d_0.pdf}\End

\Data\ID{1103107102}
\Updated{2020-12-30 10:12:38}
\Level{C}
\SubArea{Properties of functions}
\LANGen{
\Question{
The graph of the function \( f(x) \) and its inverse \( f^{-1} (x) \) is shown in one of the following pictures. Choose the picture.}
{}{}{}{}{}{}}
\LANGpl{
\Question{
Wykres funkcji \( f(x) \) i jej odwrotność \( f^{-1} (x) \) przedstawiono na rysunku poniżej. Który to rysunek?}
{}{}{}{}{}{}}
\LANGcs{
\Question{
Na jednom z níže uvedených obrázků je graf funkce \( f(x) \) a funkce \( f^{-1} (x) \) k ní inverzní. Vyberte tento obrázek.}
{}{}{}{}{}{}}
\LANGsk{
\Question{
Graf funkcie \( f(x) \) a funkcie k nej inverznej \( f^{-1} (x) \) sa nachádza na jednom z nasledujúcich obrázkov. Vyberte tento obrázok.}
{}{}{}{}{}{}}
\LANGes{
\Question{
NA}
{}{}{}{}{}{}}

\IMGanswer{1103107102a_0.pdf}
\IMGanswer{1103107102b_0.pdf}
\IMGanswer{1103107102c_0.pdf}
\IMGanswer{1103107102d_0.pdf}\End

\Data\ID{1103107103}
\Updated{2020-12-30 10:12:20}
\Level{C}
\SubArea{Properties of functions}
\LANGen{
\Question{
The graph of the function \( f(x) \) and its inverse \( f^{-1} (x) \) is shown in one of the following pictures. Choose the picture.}
{}{}{}{}{}{}}
\LANGpl{
\Question{
Wykres funkcji \( f(x) \) i jego odwrotność \( f^{-1} (x) \) przedstawione zostały na rysunku poniżej. Wybierz ten rysunek.}
{}{}{}{}{}{}}
\LANGcs{
\Question{
Na jednom z uvedených obrázků je graf funkce \( f(x) \) a funkce \( f^{-1} (x) \) k ní inverzní. Vyberte tento obrázek.}
{}{}{}{}{}{}}
\LANGsk{
\Question{
Graf funkcie \( f(x) \) a funkcie k nej inverznej \( f^{-1} (x) \) sa nachádza na jednom z nasledujúcich obrázkov. Vyberte tento obrázok.}
{}{}{}{}{}{}}
\LANGes{
\Question{
NA}
{}{}{}{}{}{}}

\IMGanswer{1103107103a.pdf}
\IMGanswer{1103107103b.pdf}
\IMGanswer{1103107103c.pdf}
\IMGanswer{1103107103d.pdf}\End

\Data\ID{1103107104}
\Updated{2020-12-30 10:12:49}
\Level{C}
\SubArea{Properties of functions}
\LANGen{
\Question{
In one of the following pictures is the graph of the function \( f \), which \textbf{is not} one-to-one function. Choose the picture.}
{}{}{}{}{}{}}
\LANGpl{
\Question{
Jeden z poniższych rysunków przedstawia wykres funkcji \( f \), gdzie \textbf{nie jest } funkcją jeden do jednego. Wybierz odpowiedni rysunek.}
{}{}{}{}{}{}}
\LANGcs{
\Question{
Na jednom z následujících obrázků je graf funkce \( f \), která \textbf{není} prostá. Vyberte tento obrázek.}
{}{}{}{}{}{}}
\LANGsk{
\Question{
Na ktorom z nasledujúcich obrázkov je graf funkcie \( f \), ktorá \textbf{nie je} prostá.}
{}{}{}{}{}{}}
\LANGes{
\Question{
NA}
{}{}{}{}{}{}}

\IMGanswer{1103107104a.pdf}
\IMGanswer{1103107104b.pdf}
\IMGanswer{1103107104c.pdf}
\IMGanswer{1103107104d.pdf}\End

\Data\ID{9000010608}
\Updated{2020-12-30 10:12:55}
\Level{C}
\SubArea{Properties of functions}
\IMG{000106_08-figure0.pdf}
\LANGen{
\Question{
Identify a function which is a possible inverse to the function graphed in the
picture.}
{\(y = x^{3}\),
\(x\in (-\infty ;\infty )\)}{\(y = x^{-3}\),
\(x\in (-2;2)\)}{\(y = x^{\frac{1} 
{3} }\),
\(x\in (0;\infty )\)}{\(y = -x^{\frac{1} 
{3} }\),
\(x\in (-\infty ;\infty )\)}{\(y = 8x\),
\(x\in (-\infty ;\infty )\)}{\(y = -4x\),
\(x\in (-\infty ;\infty )\)}}
\LANGpl{
\Question{
Z poniższej listy wybierz funkcję będącą odwrotnością funkcji, której wykres przedstawiony jest na rysunku poniżej.}
{\(y = x^{3}\),
\(x\in (-\infty ;\infty )\)}{\(y = x^{-3}\),
\(x\in (-2;2)\)}{\(y = x^{\frac{1} 
{3} }\),
\(x\in (0;\infty )\)}{\(y = -x^{\frac{1} 
{3} }\),
\(x\in (-\infty ;\infty )\)}{\(y = 8x\),
\(x\in (-\infty ;\infty )\)}{\(y = -4x\),
\(x\in (-\infty ;\infty )\)}}
\LANGcs{
\Question{
Vyberte funkci, která je inverzní k funkci, jejíž graf je na obrázku.}
{\(y = x^{3}\),
\(x\in (-\infty ;\infty )\)}{\(y = x^{-3}\),
\(x\in (-2;2)\)}{\(y = x^{\frac{1} 
{3} }\),
\(x\in (0;\infty )\)}{\(y = -x^{\frac{1} 
{3} }\),
\(x\in (-\infty ;\infty )\)}{\(y = 8x\),
\(x\in (-\infty ;\infty )\)}{\(y = -4x\),
\(x\in (-\infty ;\infty )\)}}
\LANGsk{
\Question{
Určte funkciu, ktorá je inverzná k funkcii, ktorej graf je na obrázku.}
{\(y = x^{3}\),
\(x\in (-\infty ;\infty )\)}{\(y = x^{-3}\),
\(x\in (-2;2)\)}{\(y = x^{\frac{1} 
{3} }\),
\(x\in (0;\infty )\)}{\(y = -x^{\frac{1} 
{3} }\),
\(x\in (-\infty ;\infty )\)}{\(y = 8x\),
\(x\in (-\infty ;\infty )\)}{\(y = -4x\),
\(x\in (-\infty ;\infty )\)}}
\LANGes{
\Question{
Elige la función que podría ser la inversa de la función representada por la gráfica.}
{\(y = x^{3}\),
\(x\in (-\infty ;\infty )\)}{\(y = x^{-3}\),
\(x\in (-2;2)\)}{\(y = x^{\frac{1} 
{3} }\),
\(x\in (0;\infty )\)}{\(y = -x^{\frac{1} 
{3} }\),
\(x\in (-\infty ;\infty )\)}{\(y = 8x\),
\(x\in (-\infty ;\infty )\)}{\(y = -4x\),
\(x\in (-\infty ;\infty )\)}}
\End

\Data\ID{9000010609}
\Updated{2020-12-30 10:12:24}
\Level{C}
\SubArea{Properties of functions}
\IMG{000106_09-figure0.pdf}
\LANGen{
\Question{
Identify a function which is a possible inverse to the function graphed in the
picture.}
{\(y = x^{-1}\),
\(x\in (0;\infty )\)}{\(y = x\),
\(x\in (0;\infty )\)}{\(y = -x\),
\(x\in (0;\infty )\)}{\(y = -x^{-1}\),
\(x\in (0;\infty )\)}{\(y = x^{2}\),
\(x\in (0;\infty )\)}{\(y = -x^{2}\),
\(x\in (0;\infty )\)}}
\LANGpl{
\Question{
Która z funkcji może stanowić odwrotność funkcji przedstawionej za pomocą wykresu na poniższym rysunku?}
{\(y = x^{-1}\),
\(x\in (0;\infty )\)}{\(y = x\),
\(x\in (0;\infty )\)}{\(y = -x\),
\(x\in (0;\infty )\)}{\(y = -x^{-1}\),
\(x\in (0;\infty )\)}{\(y = x^{2}\),
\(x\in (0;\infty )\)}{\(y = -x^{2}\),
\(x\in (0;\infty )\)}}
\LANGcs{
\Question{
Vyberte funkci, která je inverzní k funkci, jejíž graf je na obrázku.}
{\(y = x^{-1}\),
\(x\in (0;\infty )\)}{\(y = x\),
\(x\in (0;\infty )\)}{\(y = -x\),
\(x\in (0;\infty )\)}{\(y = -x^{-1}\),
\(x\in (0;\infty )\)}{\(y = x^{2}\),
\(x\in (0;\infty )\)}{\(y = -x^{2}\),
\(x\in (0;\infty )\)}}
\LANGsk{
\Question{
Určte funkciu, ktorá je inverzná k funkcii, ktorej graf je na obrázku.}
{\(y = x^{-1}\),
\(x\in (0;\infty )\)}{\(y = x\),
\(x\in (0;\infty )\)}{\(y = -x\),
\(x\in (0;\infty )\)}{\(y = -x^{-1}\),
\(x\in (0;\infty )\)}{\(y = x^{2}\),
\(x\in (0;\infty )\)}{\(y = -x^{2}\),
\(x\in (0;\infty )\)}}
\LANGes{
\Question{
Elige la posible inversa de la función representada en la gráfica.}
{\(y = x^{-1}\),
\(x\in (0;\infty )\)}{\(y = x\),
\(x\in (0;\infty )\)}{\(y = -x\),
\(x\in (0;\infty )\)}{\(y = -x^{-1}\),
\(x\in (0;\infty )\)}{\(y = x^{2}\),
\(x\in (0;\infty )\)}{\(y = -x^{2}\),
\(x\in (0;\infty )\)}}
\End

\Data\ID{9000010610}
\Updated{2020-12-30 10:12:57}
\Level{C}
\SubArea{Properties of functions}
\IMG{000106_10-figure0.pdf}
\LANGen{
\Question{
Identify a function which is a possible inverse to the function graphed in the
picture.}
{\(y = x^{2}\),
\(x\in (-\infty ;0] \)}{\(y = x^{-2}\),
\(x\in (-\infty ;0] \)}{\(y = -x^{2}\),
\(x\in [ 0;\infty )\)}{\(y = x^{\frac{1} 
{2} }\),
\(x\in [ 0;\infty )\)}{\(y = -x^{\frac{1} 
{2} }\),
\(x\in [ 0;\infty )\)}{\(y = -2x\),
\(x\in (-\infty ;0] \)}}
\LANGpl{
\Question{
Która z wymienionych funkcji może być odwrotnością funkcji przedstawionej za pomocą wykresu na poniższym rysunku?}
{\(y = x^{2}\),
\(x\in (-\infty ;0] \)}{\(y = x^{-2}\),
\(x\in (-\infty ;0] \)}{\(y = -x^{2}\),
\(x\in [ 0;\infty )\)}{\(y = x^{\frac{1} 
{2} }\),
\(x\in [ 0;\infty )\)}{\(y = -x^{\frac{1} 
{2} }\),
\(x\in [ 0;\infty )\)}{\(y = -2x\),
\(x\in (-\infty ;0] \)}}
\LANGcs{
\Question{
Vyberte funkci, která je inverzní k funkci, jejíž graf je na obrázku.}
{\(y = x^{2}\),
\(x\in (-\infty ;0\rangle \)}{\(y = x^{-2}\),
\(x\in (-\infty ;0\rangle \)}{\(y = -x^{2}\),
\(x\in \langle 0;\infty )\)}{\(y = x^{\frac{1} 
{2} }\),
\(x\in \langle 0;\infty )\)}{\(y = -x^{\frac{1} 
{2} }\),
\(x\in \langle 0;\infty )\)}{\(y = -2x\),
\(x\in (-\infty ;0\rangle \)}}
\LANGsk{
\Question{
Určte funkciu, ktorá je inverzná k funkcii, ktorej graf je na obrázku.}
{\(y = x^{2}\),
\(x\in (-\infty ;0\rangle \)}{\(y = x^{-2}\),
\(x\in (-\infty ;0\rangle \)}{\(y = -x^{2}\),
\(x\in \langle 0;\infty )\)}{\(y = x^{\frac{1} 
{2} }\),
\(x\in \langle 0;\infty )\)}{\(y = -x^{\frac{1} 
{2} }\),
\(x\in \langle 0;\infty )\)}{\(y = -2x\),
\(x\in (-\infty ;0\rangle \)}}
\LANGes{
\Question{
Elige la función inversa de la función representada por la gráfica.}
{\(y = x^{2}\),
\(x\in (-\infty ;0] \)}{\(y = x^{-2}\),
\(x\in (-\infty ;0] \)}{\(y = -x^{2}\),
\(x\in [ 0;\infty )\)}{\(y = x^{\frac{1} 
{2} }\),
\(x\in [ 0;\infty )\)}{\(y = -x^{\frac{1} 
{2} }\),
\(x\in [ 0;\infty )\)}{\(y = -2x\),
\(x\in (-\infty ;0] \)}}
\End
